\documentclass[11pt]{article}
\usepackage[a4paper,margin=19mm]{geometry}
\usepackage{fontspec}
\setmainfont{lmroman10-regular.otf}
\usepackage{amsmath,amssymb,amsthm,mathtools,xcolor,graphicx,hyperref}
\hypersetup{hidelinks}
\numberwithin{equation}{section}
\newtheorem{theorem}{Theorem}[section]
\newtheorem{lemma}[theorem]{Lemma}
\newtheorem{proposition}[theorem]{Proposition}
\newtheorem{corollary}[theorem]{Corollary}
\newtheorem{claim}[theorem]{Claim}
\newtheorem{fact}[theorem]{Fact}
\newtheorem{conjecture}[theorem]{Conjecture}
\theoremstyle{definition}
\newtheorem{definition}[theorem]{Definition}
\newtheorem{example}[theorem]{Example}
\newtheorem{remark}[theorem]{Remark}
\newtheorem{question}[theorem]{Question}
\theoremstyle{plain}
\newcommand{\C}{\mathbb C}
\newcommand{\R}{\mathbb R}
\newcommand{\Q}{\mathbb Q}
\newcommand{\Z}{\mathbb Z}
\newcommand{\N}{\mathbb N}
\newcommand{\eee}{\mathrm e}
\newcommand{\cclass}[1]{\textsf{#1}}
\newcommand{\cproblem}[1]{\textsc{#1}}
\newcommand{\ie}{i.\,e.}
\newcommand{\Ie}{I.\,e.}
\newcommand{\eg}{e.\,g.}
\newcommand{\Eg}{E.\,g.}
\newcommand{\phd}{Ph.\,D.}
\newcommand{\msc}{M.\,S.}
\newcommand{\bsc}{B.\,S.}
\newcommand{\expref}[2]{\hyperref[#2]{#1~\ref{#2}}}
\newcommand{\expeqref}[2]{\hyperref[#2]{#1~\eqref{#2}}}
\newcommand{\exprefsilent}[2]{\hyperref[#2]{#1}}
\providecommand{\etal}{et~al.}
\usepackage{eprint}
\usepackage{xspace}    %
\usepackage{tikz-cd}

\newtheorem{exmp}{Example}

\newcommand{\F}{\mathbb{F}}
\newcommand{\calH}{\mathcal{H}}
\newcommand{\calP}{\mathcal{P}}
\newcommand{\aff}{\mathbb{A}}
\newcommand{\proj}{\mathbb{P}}
\DeclareMathOperator{\rk}{\mathrm{rank}}
\DeclareMathOperator{\sz}{\mathrm{size}}
\DeclareMathOperator{\im}{\mathrm{Im}}
\DeclareMathOperator{\zer}{\mathrm{Zer}}
\DeclareMathOperator{\trdeg}{\mathrm{trdeg}}
\newcommand{\ol}[1]{\overline{#1}}
\newcommand{\m}{\mathfrak{m}}
\newcommand{\pder}[2]{\partial_{#2}#1}
\newcommand{\inbrace }[1]{\left\{#1\right\}} 

%
\newcommand{\Pclass}{\cclass{P}\xspace}  %
  %
%
\newcommand{\NP}{\cclass{NP}\xspace}  %
\newcommand{\coNP}{\cclass{coNP}\xspace}  %
\newcommand{\coRP}{\cclass{coRP}\xspace}  %
\newcommand{\NPcapcoNP}{\NP\,\cap\,\coNP}  %
\newcommand{\AM}{\cclass{AM}\xspace}  %
\newcommand{\coAM}{\cclass{coAM}\xspace}  %
\newcommand{\AMcapcoAM}{\AM\,\cap\,\coAM}  %
\newcommand{\SBP}{\cclass{SBP}\xspace}  %
\newcommand{\coSBP}{\cclass{coSBP}\xspace}  %
\newcommand{\SBPcapcoSBP}{\SBP\,\cap\,\coSBP}  %
\newcommand{\PSPACE}{\cclass{PSPACE}\xspace}  %
\newcommand{\EXPSPACE}{\cclass{EXPSPACE}\xspace} %
\newcommand{\VP}{\cclass{VP}\xspace}  %
\newcommand{\VPbar}{\overline{\VP}\xspace}  %
\newcommand{\PH}{\cclass{PH}\xspace}  %
\newcommand{\NC}{\cclass{NC}\xspace}  %
\newcommand{\PolyL}{\cclass{PolyL}\xspace}  %


\DeclareMathOperator{\AD}{\mathrm{AD}\xspace} %
\DeclareMathOperator{\APS}{\mathrm{APS}\xspace} %
\DeclareMathOperator{\AnnAtZero}{\mathrm{AnnAtZero}\xspace} %
\DeclareMathOperator{\HN}{\mathrm{HN}\xspace} %
\DeclareMathOperator{\NNL}{\mathrm{NNL}\xspace} %
\newtheorem{problem}[theorem]{Problem}
\title{Polynomial-space construction of hitting sets for approximable small circuits over finite fields}
\author{Zeyu Guo, Nitin Saxena, Amit Sinhababu}
\date{doi:10.4086/toc.2019.v015a016; CC BY 3.0}
\begin{document}
\maketitle
\noindent\textit{Editorial note.} Complete final author body and bibliography retained. Circuit encodings, arbitrary characteristic, original finite-field/large-q/unary-input conditions preserved. Open characteristic-zero and stronger protocol questions remain open.\par\medskip
\noindent\textit{Editorial note.} Add original ToC Problem environment interface after actual article-wrapper failure; mathematical body unchanged.\par\medskip
\noindent\textit{Allowed standard prerequisites (editorial).}\begin{itemize}
\item Affine/projective algebraic geometry, arithmetic circuits and space complexity. Original prerequisite statements/citations and full source-local fiber, projection, curve and circuit constructions retained.. No collector reproof or field/complexity strengthening. Original input encodings and finite-field scope remain explicit.
\end{itemize}\medskip
\section{Introduction}
%

Algebraic dependence is a generalization of linear dependence. Polynomials $f_1,\ldots ,f_m \in \F[x_1, \ldots ,x_n]$ are called \emph{algebraically dependent} over
the  %
field $\F$ if there exists a nonzero polynomial (called \emph{annihilator}) $A(y_1,\ldots ,y_m)\in \F[y_1,\ldots ,y_m]$ such that $A(f_1,\ldots,f_m)=0$. If no $A$ exists, then the given polynomials are called \emph{algebraically independent} over $\F$. The \emph{transcendence degree} (trdeg) of a set of polynomials is the analog of rank in linear algebra. It is defined as the
%
maximum   %
number of algebraically independent polynomials in the set. Both algebraic dependence and linear dependence
%
%
define matroid structures  %
\cite{ehrenborg1993apolarity}.
%
There exist algebraic matroids over $\F_p$ that are not
linear~\cite{ingleton1971representation}.   %

The simplest
%
example
of algebraically independent polynomials
%
is
$x_1,\ldots,x_n \in \F[x_1, \ldots$, $x_n] $. As an example of algebraically dependent polynomials, we can take $f_1=x$, $f_2=y$ and $f_3=x^2+y^2$. Then, $y_1^2 + y_2^2-y_3$ is an annihilator. The underlying field is crucial in this concept.
%
For example, for a given prime number $p$, the polynomials $x+y$ and $x^p + y^p$ are algebraically dependent over fields of characteristic $p$ (as $(x+y)^p=x^p+y^p$ in characteristic $p$).
But they are algebraically independent over fields of characteristic zero (\eg, $\mathbb{Q}$) %
and over fields of positive characteristic $\ell\neq p$.\footnote{This can be seen using the Jacobian criterion \cite{Jac41, BMS13} as discussed below.} %

%
%

Thus, the following computational question, %
%
%
$\AD(\F)$,   %
is natural and it is the first problem we consider in this paper: 
Given polynomials
%
%
$f_1,\ldots ,f_m \in \F[x_1, \ldots ,x_n]$,  %
represented as algebraic circuits, %
test if they are algebraically dependent. 
It can be solved in $\PSPACE$ using a classical result due to
Perron~\cite{Per27,Plo05}. 
We sketch the proof of this result here:
Perron proved that given a set of algebraically dependent polynomials, there exists an annihilator whose degree is
%
at most  %
the product of the degrees of the polynomials in the
%
set. (This exponential degree bound is tight~\cite{Kay09}.) %
It is not hard to see that testing the existence of such an annihilator reduces to solving a
%
system of linear equations in exponentially many variables %
where the coefficients of the annihilator are regarded as the
variables of the system.
It is also  known that solving a
%
%
%
system of linear equations %
is in logspace-uniform 
%
%
$\NC$~\cite{C76, BGH82, Mul87}, %
which is contained in $\PolyL$ (polylogarithmic space). %
%
%
Combining these results  %
we see that testing algebraic dependence
(and computing an annihilator polynomial) is in $\PSPACE$.
 
Computing the annihilator may be quite hard, but it turns out that the decision version is easy over zero (or large) characteristic using a classical result known as the Jacobian criterion~\cite{Jac41, BMS13}. The Jacobian efficiently reduces algebraic dependence testing of $f_1,\ldots,f_m$ over $\F$ to linear dependence testing of the differentials $df_1,\ldots,df_m$ over $\F(x_1,\ldots,x_n)$, where we view $df_i$ as the vector
\[
  \left( \frac{\partial f_i} {\partial x_1},\ldots, \frac{\partial f_i} {\partial x_n}\right).
\]
Placing $df_i$ as the $i$-th row gives us the Jacobian matrix $J$ of
$f_1,\ldots,f_m$. If the characteristic of the field $\F$ is zero (or
larger than the product of the degrees $\deg(f_i)$) then the
transcendence degree of $f_1,\dots,f_m$ equals $\rk(J)$.  It follows
from the
%
Polynomial Identity Lemma\footnote{Regarding the Polynomial Identity
Lemma, the following footnote appears in~\cite{CKS}. ``Variants of this
lemma, often referred to as the Schwartz--Zippel Lemma, or the
DeMillo--Lipton--Schwartz--Zippel Lemma,
were discovered at least six times, starting with
{\O}ystein Ore in 1922 and David Muller in 1954
\cite{ore,muller,schmidt,demillo1978probabilistic,zippel1979probabilistic,Sch80}.
For a brief history, see~\cite{arvind}
where the term ``Polynomial Identity Lemma''
is attributed to L. Babai.''}
that, with high probability, $\rk(J)$ is equal to the rank of $J$ evaluated at a random point in $\F^n$. This gives a simple randomized
%
polynomial-time   %
algorithm solving
%
$\AD(\F)$,   %
%
over the fields $\F$ satisfying the condition stated above.   %

For fields of positive characteristic, if the polynomials are algebraically dependent, then their Jacobian matrix is not full rank.
But the converse is not true. There are infinitely many input instances (set of algebraically independent polynomials) for which the Jacobian criterion fails. The failure can be characterized by the notion of ``inseparable extension''~\cite{pandey2016algebraic}. For example, $x^p, y^p$ are algebraically independent over $\mathbb{F}_p$, yet their Jacobian determinant vanishes. Another example is, $\left\{ x^{p-1}y, xy^{p-1} \right\}$ over $\mathbb{F}_p$ for prime $p>2$. 
%
%
Mittmann, %
Saxena, and Scheiblechner  %
\cite{mittmann2014algebraic} gave a criterion, called Witt-Jacobian,
that works over fields of
%
%
%
%
positive characteristic,
improving the complexity of
%
%
algebraic independence testing  %
from $\PSPACE$ to
%
$\NP^{\#\Pclass}$.  %
\cite{pandey2016algebraic} gave another generalization of the
Jacobian criterion that is efficient in special cases.
 
Given that an efficient algorithm to tackle 
%
positive characteristic %
is not in close sight, one could speculate the problem to be $\NP$-hard or even outside the polynomial hierarchy 
%
$\PH$. %
In
%
the present article   %
we show
%
that \emph{for finite fields,
%
$\AD(\F)$   %
is in
%
$\AMcapcoAM$}   %
(\expref{Theorem}{thm_amcoam}). This rules out the possibility of
$\NP$-hardness, %
%
unless the polynomial-time hierarchy collapses~\cite{AB09}.
%
Thus,
%
$\AD(\F)$   %
joins the league of problems of ``intermediate'' complexity,
%
%
such as   %
graph isomorphism
%
and   %
integer factoring. 
 
 \smallskip\noindent
\textbf{Constant term of the annihilators.}
We come to the second problem
\emph{AnnAtZero}  %
that we discuss in this paper: Testing if the constant term of \emph{every}
%
annihilator    %
%
of a set  %
of polynomials
%
%
(given as algebraic circuits),
%
$\mathbf{f}=\{f_1,\ldots,f_m\}$, is zero. 
The annihilators of $\mathbf{f}$ constitute a prime ideal of the polynomial ring $\F[y_1,\ldots,y_m]$. This ideal is principal when the transcendence degree of  $\mathbf{f}$ is $m-1$. This is a classical result in commutative algebra~\cite[Theorem~47]{Mat80}. See also \cite[Lemma~7]{Kay09} for an exposition. In this case, we can decide
%
in $\PSPACE$  %
if the constant term of the minimal annihilator is zero,
%
as the \emph{unique} annihilator (up to scaling) can be computed in $\PSPACE$. 
 
If
the  %
transcendence degree of $\mathbf{f}$ is less than $m-1$, the ideal of the annihilators of $\mathbf{f}$ is no longer principal. Although the ideal is finitely generated, finding the generators of this ideal is computationally very hard. (\Eg, using Gr\"obner basis techniques, we can do it in
%
$\EXPSPACE$~\cite[Section~1.2.1]{derksen2015computational}.) In this case, can we decide if all the annihilators of $\mathbf{f}$ have constant term zero? \emph{We give two equivalent characterizations of $\AnnAtZero$---one geometric and the other
%
algebraic---and we use them to  %
devise a $\PSPACE$ algorithm to solve it in all cases} (\expref{Theorem}{thm-aps}).
 
Interestingly, there is an
%
%
application of the above algorithm to algebraic complexity.  %
\emph{We give a $\PSPACE$-explicit construction of a
%
hitting set    %
of the class
%
$\VPbar_{\ol{\F}_q}$}    %
(\expref{Theorem}{thm-hsg}).
%
$\VPbar_{\ol{\F}_q}$    %
consists of $n$-variate
%
%
polynomials of degree $d=n^{O(1)}$ over
the field $\ol{\F}_q$, that can be ``infinitesimally approximated'' by
%
%
algebraic circuits of size $s=n^{O(1)}$.
A polynomial $p(x_1,\ldots,x_n)$ over an algebraically closed field $\ol{\F}$ is said to be \emph{infinitesimally approximated} by a circuit of size $s$, if the circuit computes a polynomial of the form $p+\epsilon p_1+\epsilon^2 p_2+\cdots+\epsilon^m p_m$, where the circuit uses constants from $\ol{\F}(\epsilon)$ (for example, ${1}/{\epsilon}$ can be used as a constant) and $p_1,\ldots,p_m \in \ol{\F}[x_1,\ldots,x_n]$.

%
A set $\mathcal{H}$ of points   %
is called a
%
\emph{hitting set}    %
for a
%
set $C$ of polynomials,   %
if any nonzero polynomial in $C$ evaluates to a nonzero value in at least one point in $\mathcal{H}$. 
It was shown by  %
Heintz and Schnorr   %
\cite{heintz1980testing} that
%
hitting sets  %
of size $\mathrm{poly}(n, s)$ and bit length $\mathrm{poly}(n, s,\log d)$
\emph{exist} for the class of $n$-variate polynomials of degree
at most $d$ computed by size $s$ arithmetic circuits.
The same result extends to  
the class of $n$-variate polynomials of degree at most $d$ \emph{infinitesimally approximated} by size $s$ arithmetic circuits. Now the
%
%
problem of constructing explicit hitting sets  %
for
%
$\VPbar$    %
asks to
%
%
%
deterministically compute a   %
%
set of points that is a hitting set for the class of $n$-variate polynomials of degree at most $d$, infinitesimally approximated by size $s$ arithmetic circuits.
%

The above problem is interesting as natural questions like explicit construction of the normalization map (in Noether's Normalization Lemma $\NNL$) reduce to the construction of a hitting set of
%
$\VPbar$    %
%
(Mulmuley~\cite{mulmuley2017geometric}), %
%
which was previously known to be only in
%
$\EXPSPACE$~\cite{mulmuley2017geometric,mul12}. 
This was recently put in $\PSPACE$, over the field $\mathbb{C}$,
by  %
Forbes and Shpilka~\cite{forbes2017pspace}.
Their proof
%
uses real analysis and does not apply to finite fields. We need to develop purely algebraic concepts to solve the finite field case (namely through $\AnnAtZero$), which apply to \emph{any} field.
%
%
%
%

\smallskip 
To further motivate the concept of algebraic dependence, we list a few recent problems in computer science.
 The first problem is about constructing an explicit randomness extractor for sources which are polynomial maps over finite fields. Using the Jacobian criterion, \cite{DGW09, bib:Dvi09} solved the problem for fields with large characteristic.
 The second application is
%
to   %
the famous polynomial identity testing (PIT) problem. To efficiently design hitting sets, for some interesting models, 
 \cite{BMS13, ASSS11, KS15} constructed a family of trdeg-preserving maps. For more background and applications of algebraic dependence testing, see~\cite{pandey2016algebraic}. The annihilator has been a key concept to prove the connection between hitting sets and lower bounds~\cite{heintz1980testing}, and in  bootstrapping ``weak'' hitting sets~\cite{AGS17}. 

Finally, we remark that a common thread among the problems we study in this paper is studying the closure of the polynomial map.
Let $f:\F^n \to \F^n$ be a given polynomial map. Then, we study two computational questions: (1) (Algebraic Dependence Testing) whether 
the %
dimension of 
%
the Zariski closure of the image of $f$  %
is $< n$? (2) (Origin in closure) whether 
%
the origin $\mathbf{0}$ is in the Zariski closure of the image of $f$? %

%
\subsection{Our results}\label{contrib}
%

In this paper, we give Arthur-Merlin protocols
%
and   %
algorithms, with proofs using basic tools from algebraic geometry. The first theorem we prove is about
%
$\AD(\F_q)$\,.   %

\begin{theorem}\label{thm_amcoam}
%
%
%
Testing algebraic independence of multivariate polynomials over a
finite field is in
%
$\AMcapcoAM$.
%
%
\end{theorem}
 
This result vastly improves the current best upper bound known
for
%
$\AD(\F_q)$---from  %
being ``outside'' the polynomial hierarchy (namely
%
$\NP^{\#\Pclass}$   %
\cite{mittmann2014algebraic}) to ``lower'' than the
%
second level   %
of
the  %
polynomial hierarchy (namely
%
$\AMcapcoAM$).    %
This rules out the possibility
%
of $\AD(\F_q)$   %
being $\NP$-hard (unless
the  %
polynomial hierarchy collapses to the
%
second level   %
\cite{AB09}).  %
Recall that,
%
for a field $\F$ of characteristic zero or of large characteristic,  %
%
$\AD(\F)$   %
is in
%
$\coRP$  %
(\expref{Section}{sec-prelim}). We conjecture such a result for
%
$\AD(\F_q)$    %
too.

\smallskip
Our second result is about the problem $\AnnAtZero$ (\ie, testing whether all the annihilators of given $\mathbf{f}$ have constant term zero). A priori it is unclear why it should  have complexity better than
%
$\EXPSPACE$.  %
%
(Note: ideal    %
membership is
%
$\EXPSPACE$-complete~\cite{mayr1982complexity}.)   %
%
%
%
%
%
%


First,   %
we relate %
the problem $\AnnAtZero$  %
to a (new) version of polynomial system satisfiability,
over the algebraic closure $\overline{\F}$. %
%


%
%
\begin{problem}[Approximate polynomials satisfiability ($\APS$)]
Given polynomials
%
%
\[
  f_1,\ldots,f_m \in \F[x_1,\ldots,x_n]\,,
\]
represented as algebraic circuits, %
does there exist $ \mathbf\beta \in \overline{\F}(\epsilon)^n $ such
that for all $i$, $f_i(\mathbf{\beta})$ is in the ideal
$\epsilon \overline{\F}[\epsilon]$ of $\overline{\F}[\epsilon]$? If
yes, then we say that $\mathbf{f}:= \{f_1,\ldots,f_m\}$ is in $\APS$.
\end{problem}

%
It is easy to show that the function %
field $\ol{\F}(\epsilon)$ here can be equivalently replaced by the ring $\overline{\F}[\epsilon, \epsilon^{-1}]$  of \emph{Laurent polynomials}, or, the field $\ol{\F}((\epsilon))$ of \emph{formal Laurent series} (use mod $\epsilon \overline{\F}[\epsilon]$). A reason why these objects appear in algebraic complexity can be found in~\cite[Section~5.2]{burgisser2004complexity} and~\cite[Section~5]{LL89}. They help algebrize the notion of ``infinitesimal approximation'' (in real analysis think of $\epsilon\rightarrow0$ \& $1/\epsilon\rightarrow\infty$). A notable computational issue involved is that the degree bound of $\epsilon$ required for $\beta$ is exponential in the input size~\cite[Proposition~3]{LL89}; this may again be a ``justification''
%
why $\APS$ may require  %
that much space.

\smallskip
Classically, the \emph{exact} version of $\APS$ has been
%
well-studied---Does there exist $\mathbf{\beta} \in \overline{\F}^n$ such that for all $i$, $f_i(\mathbf{\beta})=0 $? This is what Hilbert's Nullstellensatz 
%
($\HN$) %
%
characterizes. 
%
$\HN$ also yields %
an impressive $\PSPACE$ algorithm~\cite{koiran1996hilbert, kollar1988sharp}. Note that if system $\mathbf{f}$ has an exact solution, then it is trivially in $\APS$. But the converse is not true. For example,  $\{x, xy-1\}$ is in $\APS$, but there is no exact solution in $\overline{\F}$. To see that it is in $\APS$, assign $x=\epsilon$ and $y= 1/\epsilon$. Also, the instance $\{x, x+1\}$ is neither in $\APS$ nor
%
does it have  %
an exact solution. Finally, note that if we restrict $ \mathbf{\beta}$ to come from $\overline{\F}[\epsilon]^n$ then $\APS$ becomes equivalent to exact satisfiability and
%
$\HN$ %
applies. This can be seen by going modulo $\epsilon \overline{\F}[\epsilon]$, as the quotient $\overline{\F}[\epsilon]/\epsilon \overline{\F}[\epsilon]$ is $\overline{\F}$.

Coming back to $\AnnAtZero$, we show that it is equivalent both to a geometric question and to deciding $\APS$. This gives us, with more work, the following surprising consequence.

\begin{theorem}\label{thm-aps}
 $\APS$ is $\NP$-hard and is in $\PSPACE$.
\end{theorem}

We apply this to designing hitting sets and solving
%
$\NNL$.  (We refer to~\cite{mulmuley2017geometric} for the background.)


\begin{theorem}\label{thm-hsg}
  There is a $\PSPACE$ algorithm that 
 %
  (given input $n,s,r$ in unary
%
and  %
suitably large $q$ in binary)
%
computes  %
a set of points from $\F_q ^n$ of size 
%
$\mathrm{poly}(n,s)$ %
%
that hits all $n$-variate degree-$r$ polynomials over $\ol{\F}_q $ that can be infinitesimally approximated by
%
size-$s$     %
circuits.
\end{theorem}

%

To state the results on $\NNL$, we need the following
%
notation    %
from~\cite{mulmuley2017geometric}. %
Let $m$ be a positive integer and $r=m^2$.   Let $V$ be the vector space of homogeneous degree-$m$ polynomials in the variables $x_1,\dots, x_{r}$ over an algebraically closed field $\F$. 
Each $f\in V$ determines a point $[f]$ in the projective space $\proj V$.
Let $\Sigma[\mathrm{det}, m]$ be the set of points $[f]$  such that $f$ is a symbolic determinant of a matrix $M\in \mathbb{F}^{m\times m}$ whose entries are linear forms in $x_1,\dots,x_r$. 
Finally, let $\Delta[\mathrm{det}, m]\subseteq \proj V$ be the
%
Zariski   %
closure of $\Sigma[\mathrm{det}, m]$. 

%
Mulmuley~\cite{mulmuley2017geometric} considers the problem of constructing a homogeneous linear map that maps $\Sigma[\mathrm{det}, m]\subseteq\proj V$ to a much smaller projective space $\proj^k$ (a \emph{normalizing map}). In particular, we want to compute a set $S=\{p_0,\dots,p_k\}$ in $\mathrm{poly}(m)$-time, where $k= \mathrm{poly}(m)$ and $p_i\in \F^r$ for $i\in [k]$, such that the map $f\mapsto (f(p_0),\dots,f(p_k))$ induces a
%
well-defined    %
map from $\Sigma[\mathrm{det}, m]$ to $\proj^k$. Such a set $S$ is called a (strict) \emph{e.s.o.p.}\ (\emph{explicit system of parameters}). In addition, an e.s.o.p.\ $S$ is \emph{separating} if for any distinct $f,g$
%
%
satisfying $[f], [g]\in \Delta[\mathrm{det}, m]$,
we have $f(p_i)\neq g(p_i)$ for some $p_i\in S$.

As shown in~\cite{mulmuley2017geometric}, the problem of constructing a
separating e.s.o.p.\ reduces to that of constructing hitting sets for 
%
$\VPbar$ %
\cite[Theorem~4.5]{mulmuley2017geometric}. Combining this reduction with our $\PSPACE$ algorithm of constructing hitting sets for 
%
$\VPbar$, %
we obtain the following result.
\begin{theorem}[{$\NNL$ for $\Delta[\mathrm{det}, m]$ in $\PSPACE$}] 
The problem of constructing a separating e.s.o.p.\  for $\Delta[\mathrm{det}, m]$ belongs to $\PSPACE$.
\end{theorem}

The results for $\Delta[\mathrm{det}, m]$ also hold for other  varieties, as long as such varieties satisfy some explicitness condition. This is captured by the notion of an \emph{explicit family $\{W_n\}$ of varieties} in~\cite{mulmuley2017geometric}. See~\cite[Definition~5.1]{mulmuley2017geometric} for the formal definition. For such varieties $W_n$, constructing a separating e.s.o.p.\ also reduces to constructing hitting sets for 
%
$\VPbar$ %
\cite[Theorem~5.11]{mulmuley2017geometric}.
%
So we have the following result.
\begin{theorem}[{$\NNL$ for explicit varieties in $\PSPACE$}] 
Let $\{W_n\}$ be an explicit family of  varieties as in~\cite[Definition~5.1]{mulmuley2017geometric}.
Then the problem of constructing a separating e.s.o.p.\  for $W_n$  belongs to $\PSPACE$.
\end{theorem}

%

\smallskip
\noindent
\textbf{More applications?} The exact polynomials satisfiability question 
%
$\HN$ %
(over $\ol{\F}$) is highly expressive and, naturally,
%
many computational problems
%
%
%
%
can be  %
expressed that way. We claim that in a similar spirit, the $\APS$ question expresses those computer science problems
%
that involve ``infinitesimal approximation.''
%
%
%
%
Since finite fields do not have a topology  %
allowing approximations, algebraic approximations over arbitrary fields
%
are   %
needed. The latter has been useful in fast matrix multiplication algorithms. 


One prominent example of algebraic approximation is the concept of \emph{border rank} of tensor polynomials (used in matrix multiplication algorithms and
%
Geometric Complexity Theory (GCT),  %
see~\cite{burgisser2013algebraic, landsberg2012tensors, le2014powers}). Border rank computation of a given tensor (over $\ol{\F}$) can easily be reduced to an $\APS$ instance and, hence, now solved in $\PSPACE$. This brings the complexity of border rank closer to the complexity of tensor rank~\cite{schaefer2016complexity}. {}From the point of view of Gr\"obner basis theory, $\APS$ is a problem that seems a priori much harder than 
%
$\HN$. %
Now that both of them have a $\PSPACE$ algorithm, one may wonder whether it can be brought all the way down to
$\NP$ or $\AM$?  %
%
(In fact, $\HN_{\mathbb{C}}$  %
is known to be in $\AM$, conditionally under GRH~\cite{koiran1996hilbert}.)

More generally, $\APS$ is naturally related to GCT as
%
follows. The   %
program of GCT aims to find an explicit hard polynomial $g$ (\eg, the permanent) in a vector space $V\subseteq\mathbb{F}[x_1,\dots,x_n]$ and prove that its projective image $[g]\in\proj V$ is not contained in a certain projective variety $X$ that consists of  the ``easy'' polynomials.
The \emph{affine cone} $\hat{X}\subseteq V$ of $X$ is often specified as the Zariski closure of the image of some morphism $\mathbf{f}$  from an affine space (see, \eg, \cite[Section~2]{Rai12} for the case that $X$ is a higher secant variety of a Segre--Veronese variety). The problem then becomes proving $g\not\in \hat{X}=\overline{\im(\mathbf{f})}$, or equivalently, proving $\mathbf{0}\not\in \overline{\im(\mathbf{f}')}$, where $\mathbf{0}$ is the origin of $V$ and $\mathbf{f}'$ is the translation of $\mathbf{f}$ by $-g$. As we will see in this paper, one equivalent formulation of $\APS$ is precisely the problem of deciding if $\mathbf{0}$ is in the Zariski closure of the image of a given morphism from an affine space. 


Our methods in the proof of \expref{Theorem}{thm-aps} also imply an interesting ``degree bound'' related to the (prime) ideal $I$ of annihilators of
%
the set $\mathbf f$ of polynomials.  %
Namely, $I=\sqrt{I_{\le d}}$, where $I_{\le d}$ refers to the subideal generated by degree $\le d$ polynomials
%
in    %
$I$, $d$ is the Perron-like bound $(\max_{i\in [m]}\deg(f_i))^k$, and $k:= \trdeg(\mathbf f)$. This is equivalent to the geometric fact, which we prove, that the varieties defined by the two ideals $I$ and $I_{\le d}$ are equal (\expref{Theorem}{thm-randReduct}). This again is an exponential improvement over what one expects to get from the general Gr\"obner basis
methods, %
because, the generators of $I$ may well have
%
doubly exponential   %
degree.

The result on hitting set (\expref{Theorem}{thm-hsg}) can be applied to compute, in $\PSPACE$, the explicit system of parameters (e.s.o.p.) of the \emph{invariant ring} of the variety $\Delta[\det,s]$, over $\ol{\F}_q$, with a given group action~\cite[Theorem~4.9]{mulmuley2017geometric}. Also, we can now construct, in $\PSPACE$, polynomials in $\F_q[x_1,\dots,x_n]$ that cannot even be approximated by ``small'' algebraic circuits. Such results were previously known only for
%
fields of characteristic zero,  %
see~\cite[Theorems~1.1--1.4]{forbes2017pspace}. Bringing this complexity down to 
%
$\Pclass$  %
is the longstanding problem of blackbox PIT
%
(and lower bounds),  %
see~\cite{Saxena09, shpilka2010arithmetic, Saxena14}. Mulmuley~\cite{mul12} pointed out that small hitting sets for 
%
$\VPbar$ %
can be designed in
%
$\EXPSPACE$
which is a far worse complexity than that for 
%
$\VP$. %
He called it the GCT Chasm. We bridge it somewhat, as the proof of \expref{Theorem}{thm-hsg} shows that small hitting sets for 
%
$\VPbar_{\ol{\F}}$ %
can be designed in $\PSPACE$ (like those for 
%
$\VP$) %
for \emph{any} field $\F$.


In another application, the null-cone problem defined in~\cite{brgisser_et_al:LIPIcs:2018:8351} can be seen as a special case of $\APS$ and using our algorithm, it can be solved in $\PSPACE$. B\"urgisser et al.~\cite{brgisser_et_al:LIPIcs:2018:8351} gave an 
%
exponential-time  %
algorithm for the above problem (bringing it down from
%
$\EXPSPACE$).   %
%

Finally, another motivation for $\AnnAtZero$ or $\APS$ comes from the study of
the  %
\emph{geometric ideal proof system} in algebraic proof complexity, introduced in~\cite{grochow2018circuit}, where they essentially discuss $\AnnAtZero$ for systems of polynomial equations corresponding to Boolean tautologies.
See \cite[Appendix B]{grochow2018circuit} for more details.


%

\smallskip
\noindent
%
%
\begin{remark}
Given a system of polynomial equations $f_1(x_1, \ldots, x_n)=0,\ldots,f_m(x_1, \ldots, x_n)=0$, we can always homogenize
%
the $f_i$  %
by introducing a new variable $z$. The question of \emph{projective polynomials satisfiability} for a given system of equations is whether there is a nonzero solution  (called a \emph{projective solution}) to the homogenized system. We have seen that $\APS$ does not directly reduce to (affine) polynomials satisfiability, as there are unsatisfiable systems, \eg, $\{X=0,~XY=1\}$, that have approximate solutions. We note that $\APS$ does not directly reduce to projective polynomials satisfiability
%
either.  Consider, for example,  %
the system $\{X+Y=0,\;X+Y=1\}$, corresponding to two parallel affine lines. It has no approximate solution, but the homogenized system $\{X+Y=0,\; X+Y=Z\}$ has a projective solution $(1, -1, 0)$. 

Nonetheless, it is true that if the equations $f_1=0, \dots, f_m=0$ have an approximate solution, then the homogenized equations $\hat{f}_1=0,\dots,\hat{f}_m=0$  have a projective solution (in $\proj^n$). We sketch a proof
%
of    %
this fact. %
Suppose $a_1,\dots,a_n \in \overline{\F}(\epsilon)\subseteq  \overline{\F}((\epsilon))$ form an approximate solution of the original system.
%
 Let $a_{n+1}=1$.
Then   $f_i(a_1,\dots,a_n)=\hat{f}_{i}(a_1,\dots,a_n,a_{n+1})$ for $i\in [m]$.
Choose the smallest $k\in\Z$ such that $\epsilon^k a_i$ is in the ring of formal power series $\overline{\F}[[\epsilon]]$ for all $i\in [n+1]$.
We have $k\geq 0$ as $a_{n+1}=1$. 
For $i\in [n+1]$, let $\bar{a}_i\in \overline{\F}$ be the constant term of $\epsilon^k a_i$.
Minimality of $k$ guarantees that   $\bar{a}_i\neq 0$  for some $i\in [n+1]$.
Assigning $\bar{a}_1,\dots,\bar{a}_{n+1}$ to $x_1,\dots,x_n,z$ then gives a projective solution to the equations $\hat{f}_1=0,\dots,\hat{f}_m=0$.
\end{remark}  %

%
%
%


 %
 \subsection{Proof ideas}\label{idea}
%
 
\smallskip\noindent
\textbf{Proof idea of \expref{Theorem}{thm_amcoam}.}  
Suppose we are given polynomials (represented as algebraic circuits) $\mathbf{f}:= \{f_1,\ldots,f_m\}$ computing in $\F_q[x_1,\dots,x_n]$. For the
$\AM$ and $\coAM$ protocols, we consider the following system of equations over a ``small'' extension $\F_{q'}$. %
%
  
 For $b=(b_1,\dots,b_m)\in \F_{q'}^m$, define the system of equations $f_i(x_1,\dots,x_n)=b_i$, for $i\in [m]$. We denote the number of solutions of the above system in $\F_{q'}^n$ as $N_b$. Let $f:\F_{q'}^n\to \F_{q'}^m$ be the polynomial map $a\mapsto (f_1(a),\dots,f_m(a))$.

 \emph{$\AM$ gap.} [\expref{Theorem}{thm_am}] We establish bounds for the number $N_{f(a)}$, where $a$ is a random point in $\F_{q'}^n$. If $f_1,\dots,f_m$ are independent, we show that  $N_{f(a)}$ is relatively small.
%
On the other hand,   %
if the polynomials are algebraically dependent then $N_{f(a)}$ is much
%
greater.   %

Assume $\mathbf{f}$
%
is   %
algebraically independent. Without loss of generality~\cite[Section~2]{pandey2016algebraic} we can assume that $m=n$ and for all $i\in[n]$, 
%
$x_i,f_1,\ldots,f_n$ %
are algebraically dependent. 
The first step is to show that
%
 with high probability, the zero set defined by the system of equations, for random $f(a)$, is finite (\ie, it has dimension $\leq 0$ as a variety). %
%
This is proved using the Perron degree bound on the annihilator of $\{x_i,f_1,\ldots,f_n\}$. 
Next, one can apply an affine version of 
%
B\'ezout's %
theorem to
%
obtain an upper bound on   %
$N_{f(a)}$. 
On the other hand, suppose $\mathbf{f}$
%
is  %
algebraically dependent, say with annihilator $Q$. Let $\im(f):=f(\F_{q'}^n)$ be the image of $f$. Since $Q$ vanishes on $\im(f)$, we know that $\im(f)$ is relatively small, whence we deduce that $N_{f(a)}$ is large for ``most''
%
values of $a$.   %
 
\emph{$\coAM$ gap.} [\expref{Theorem}{thm_coam}] We pick a random point $b=(b_1,\dots,b_m)\in \F_{q'}^m$ and bound $N_b$, which is the number of solutions of the system defined above.
In the dependent case, we show that $N_b=0$ for ``most''
%
values of $b$.  %
But in the independent case, we show that $N_b\ge1$ for ``many''
%
(maybe   %
not ``most''!)
%
values of $b$.   %
The ideas are based on those sketched above.
%

The two kinds of gaps shown above are based on the
%
sets   %
%
$f^{-1}(f(\mathbf x))$
%
and $\im(f)$, resp.   %

Note that membership in either of these sets is testable in $\NP$ (the latter requires nondeterminism). Based on this and the gaps between the respective cardinalities, we can invoke \expref{Lemma}{lem_amprotocol} and devise the
$\AM$ and $\coAM$ protocols for
%
$\AD(\F_{q'})$,    %
which also apply to
%
$\AD(\F_q)$.    %

%
\begin{remark}   %
%
Our proof of \expref{Theorem}{thm_amcoam} employs the fact that we could efficiently sample %
a random point in the set $\im(f)$. In contrast, it is not clear how to 
efficiently %
sample a random point in the zero set $\zer(\mathbf f):= \{\mathbf x\in \F_{q'}^n \,\mid\, f(\mathbf x)= \mathbf 0\}$. Thus, we manage to side-step the $\NP$-hardness associated with most zero set properties. \Eg, computing the dimension of $\zer(\mathbf f)$ is $\NP$-hard.
\end{remark}
  
\smallskip\noindent
\textbf{Proof idea of \expref{Theorem} {thm-aps}.} 
Let polynomials 
%
$\mathbf f:= \{f_1,\dots,f_m\}$ in $\F[x_1,\dots,x_n]$ be given over a field $\F$,
represented by algebraic circuits. %
We want to determine if the constant term of every annihilator for $\mathbf f$ is zero. Redefine the polynomial map $f:\ol{\F}^n\to \ol{\F}^m$; $a\mapsto (f_1(a),\dots,f_m(a))$. For a subset $S$ of an affine
%
or projective   %
space, write $\overline{S}$ for its \emph{Zariski closure} in that space, \ie,
%
for  %
the smallest subset that contains $S$ and equals the zero set $\zer(I)$ of some polynomial %
ideal $I$ (homogeneous ideal in the projective case).  %

\emph{$\APS$ vs.\ $\AnnAtZero$.} [\expref{Theorem}{thm_approx}] Now, we interpret the problem $\AnnAtZero$ in a geometric way through \expref{Lemma}{lem_geom}:

 The constant term of every annihilator of $\mathbf f$ is zero iff the origin point $\mathbf 0\in \ol{\im(f)}$. 
 
This has a simple proof using the ideal-variety correspondence~\cite{Har92}. Note that the stronger condition $\mathbf 0\in\im(f)$ is equivalent to the existence of a common solution to the equations $f_i(x_1,\dots,x_n)$ $=0$, $i=1,\dots,m$. The latter problem (call it 
%
$\HN$ %
for Hilbert's Nullstellensatz) is known to be in $\AM$ if  $\F=\mathbb{Q}$ and GRH is assumed~\cite{koiran1996hilbert}.
However, $\im(f)$ is not necessarily Zariski closed; equivalently, it may be strictly smaller than $\ol{\im(f)}$. So, we need new ideas to test $\mathbf{0}\in \overline{\im(f)}$. 
  
Next, we observe that although $\mathbf{0}\in \overline{\im(f)}$ is not equivalent to the existence of a solution $\mathbf x\in \ol{\F}^n$ to $f(\mathbf x)= \mathbf 0$, it \emph{is} equivalent to the existence of an  ``approximate solution'' $\mathbf x\in \ol{\F}(\epsilon)^n$, which is an $n$-tuple of rational functions in a formal variable $\epsilon$. The proof idea of this uses a degree bound on $\epsilon$ due to~\cite{LL89}. We called this problem $\APS$. As $\AnnAtZero$ problem is already known to be $\NP$-hard~\cite{Kay09}, $\APS$ is also $\NP$-hard. 

%
\emph{Upper bound on $\APS$.}   %
We now know 
%
that solving %
 $\APS$ for $\mathbf f$ is equivalent to solving $\AnnAtZero$ for $\mathbf f$. $\AnnAtZero$ was previously known to be in $\PSPACE$ in the special case when the transcendence degree $k$ of $\F(\mathbf f)/\F$ equals $m$ or $m-1$, but the general case remained open (best being
%
$\EXPSPACE$).   %

In this
%
article  %
we prove that $\AnnAtZero$ is in $\PSPACE$ even when $k<m-1$. Our simple idea is to reduce the input to a smaller $m=k+1$ instance, by choosing new polynomials $g_1,\dots,g_{k+1}$ that are random linear combinations of 
%
the polynomials $f_i$. %
We show that with high probability, replacing $\{f_1,\dots,f_m\}$ by $\{g_1,\dots,g_{k+1}\}$ preserves YES/NO instances as well as the transcendence degree. This gives a randomized poly-time reduction from the case $k<m-1$ to $k=m-1$ (\expref{Theorem}{thm-randReduct}). The latter has a standard $\PSPACE$ algorithm.

For notational convenience view $\ol{\F}$ as the \emph{affine line} $\aff$.
Define $V:=\overline{\im(f)}\subseteq \aff^m$. Proving that the above reduction (of $m$) does preserve YES/NO instances amounts to proving the following geometric statement: If $V$ does not contain the origin $O\in\aff^m$, then with high probability, the variety $V':=\overline{\pi(V)}$ does not contain the origin $O'\in\aff^{k+1}$ either, where $\pi:\aff^m\to \aff^{k+1}$ is a random linear map. 

As $\pi$ is picked at random, the kernel $W$ of $\pi$ is a random linear subspace of $\aff^m$.
We have $O'\not\in \pi(V)$ whenever $V\cap  W=\emptyset$, but this is not sufficient for proving $O'\not\in \overline{\pi(V)}$, since $V$ may ``get arbitrarily close to $W$'' in $\aff^m$ and meet $W$ ``at infinity.'' 

\begin{exmp}
The hyperbola $V:=V(X_1 X_2-1)$ and the line $W:=V(X_1)$ do not meet in the affine space $\aff^2$, even though they get closer and closer as we increase the absolute value of the coordinate $X_2$.
However, note that the  projective closures of these two varieties $V_c:=V(X_1X_2-X_0^2)$ and $W_c:=V(X_1)$ do meet at the point $(0,0,1)$ of  the projective space $\proj^2$ , which is thought of as one of the ``points at infinity.'' %
\end{exmp}

Inspired by the above observation, we consider projective geometry instead of affine geometry, and prove that $O'\not\in V'$ holds as long as the projective closure of $V$ and that of $W$ are disjoint. The proof uses a construction of a projective subvariety---the \emph{join}---to characterize $\pi^{-1}(V')$, and eventually rules out $W\subseteq \pi^{-1}(V')$ (\expref{Lemma}{lem_projcriterion2}).

Moreover, we show that this holds with high probability if 
%
$O\not\in V$, by (repeatedly) using the fact that a general (=random) %
hyperplane section reduces the dimension of a variety by one.

\smallskip\noindent
\textbf{Proof idea of \expref{Theorem}{thm-hsg}.} 
Define $\aff:=\ol{\F}_q$ and assume w.l.o.g.\ $q\ge\Omega(sr^{2})$~\cite{AL86}.
\cite[Theorem~4.4]{heintz1980testing} showed that a hitting set, of size $h:= O(s^2n^2)$ in $\F_q^n$, \emph{exists} for the class of degree-$r$ polynomials, in $\aff[x_1,\dots,x_n]$, that can be infinitesimally approximated by size-$s$ algebraic circuits. So, we can search over all possible subsets of size $h$ from $\F_q^n$ and ``most'' of them are hitting sets.

How do we certify that a candidate set $\calH$ is a hitting set? The idea is to use universal circuits. A \emph{universal circuit} has $n$ essential variables $\mathbf x=\{x_1,\ldots,x_n\}$ and $s':=O(sr^4)$ auxiliary variables $\mathbf y= \{y_1,\ldots, y_{s'}\}$. We can fix the auxiliary variables, from $\aff(\epsilon)$, in such a way so that it can output any homogeneous circuit of size-$s$, approximating a degree-$r$ polynomial in
%
$\VPbar_\aff$.   %
Given a universal circuit $\Psi$, certification of a hitting set $\calH$ is based on the following observation, that follows from the definitions: 

%
A candidate set %
$\calH=: \{\mathbf v_1,\ldots, \mathbf v_h\}$ is a hitting set iff
\[
  \forall \mathbf y\in \aff(\epsilon)^{s'}, \,\Psi(\mathbf{y,x})\notin \epsilon\aff[\epsilon][\mathbf x]  \,\Rightarrow\, \exists i\in[h],\,  \Psi(\mathbf y, \mathbf v_i)\notin \epsilon\aff[\epsilon].
\]

\smallskip
%
Equivalently, a candidate set %
$\calH= \{\mathbf v_1,\ldots, \mathbf v_h\}$ is \emph{not} a hitting set iff
\[
  \exists \mathbf y\in \aff(\epsilon)^{s'}, \,\Psi(\mathbf{y,x})\notin \epsilon\aff[\epsilon][\mathbf x]\,\qquad\text{and}\qquad  \forall i\in[h],\,  \Psi(\mathbf y, \mathbf v_i)\in \epsilon\aff[\epsilon]\,.
\]

\smallskip
Note that certification of hitting set is more challenging than the one against polynomials in
%
$\VP$,  %
because the degree bounds for $\epsilon$ are exponentially high and moreover, we do not know how to frame the first ``non-containment'' condition as an $\APS$ instance. To translate it to an $\APS$ instance, our key idea is the following. 

Pick $q\ge\Omega(s'r^2)$ so that a hitting set exists, in $\F_q^n$, that works against polynomials infinitesimally approximated by the specializations of $\Psi$. 
Suppose $\Psi(\mathbf{\alpha,x})$ is not in $\epsilon\aff[\epsilon][\mathbf x]$, for some $\mathbf \alpha\in \aff(\epsilon)^{s'}$. This means that we can write it as $\sum_{-m\le i\le m'}$ $\epsilon^i g_i(\mathbf x)$ with $g_{-m}\ne0$ and $m\ge0$. Clearly, $\epsilon^m\cdot \Psi(\mathbf{\alpha,x})$ infinitesimally approximates the nonzero polynomial $g_{-m} \in \aff[\mathbf x]$.
By the conditions on $\Psi$, we know that $g_{-m}$ is a homogeneous 
%
polynomial of degree $r$ %
(and approximative complexity $s'$). Thus, by 
%
%
%
\expref{Lemma}{Polynomial Identity},  %
there exists a $\mathbf \beta\in\F_q^n$ such that $g_{-m}(\mathbf \beta)=:a$ is a nonzero element in $\aff$. We can normalize by this and consider $a^{-1}\epsilon^m\cdot \Psi(\mathbf{y,x})$, which evaluates to $1+\epsilon\aff[\epsilon]$ at $(\mathbf{\alpha,\beta})$. Since this normalization factor only affects the auxiliary variables $\mathbf y$, we get another equivalent criterion:

%
A candidate set %
$\calH= \{\mathbf v_1,\ldots, \mathbf v_h\}$ is \emph{not} a hitting set iff 
$\exists \mathbf y\in \aff(\epsilon)^{s'}$ and $\exists \mathbf x\in \F_q^n$  such that,
  \[
    \Psi(\mathbf{y,x}) - 1 \in \epsilon\aff[\epsilon]\quad\text{and}\quad \forall i\in[h],  \Psi(\mathbf y, \mathbf v_i)\in \epsilon\aff[\epsilon].
  \]

\smallskip
We reach closer to $\APS$, but how do we implement $\exists? \mathbf x\in \F_q^n $ (it takes exponential space)? 

The idea is to rewrite it, instead using the $(r+1)$-th roots of unity $Z_{r+1}\subseteq\aff$, 
%
as %
$\exists \mathbf x\in \aff(\epsilon)^n$, $\forall i\in[n]$, $x_i^{r+1}-1 \in \epsilon\aff[\epsilon]$. This gives us a criterion that is an instance of $\APS$ with $n+h+1$ input polynomials (\expref{Theorem}{thm_HITTING}). By \expref{Theorem}{thm-aps} it can be done in 
%
$\PSPACE$, %
finishing the proof. Moreover, this $\PSPACE$ algorithm 
%
is independent of 
%
the characteristic of the underlying field. %
(\Eg, it can be seen as an alternative to~\cite{forbes2017pspace} over the complex field.)

%

\section{Preliminaries}\label{sec-prelim}
%

\textbf{Jacobian.}
Although this work would not need it, we define the classical Jacobian: 
%
Given polynomials $\mathbf{f} = \inbrace{f_1, \cdots, f_m}$ in $\F[x_1,\cdots, x_n]$, the \emph{Jacobian} of $\mathbf{f}$ %
is the matrix $\mathcal{J}_{\mathbf{x}}(\mathbf{f}) := (\pder{f_i}{x_j})_{m \times n}$, where $\pder{f_i}{x_j} := \partial f_i/ \partial x_j$. 

The Jacobian criterion~\cite{Jac41, BMS13}
%
states that for  %
degree $\le d$ and trdeg $\le r$ polynomials $\mathbf{f}$, if char$(\F) = 0$ or char$(\F) > d^r$, 
then $\trdeg(\mathbf{f}) \,=\, \rk_{\F(\mathbf{x})} \mathcal{J}_{\mathbf{x}}(\mathbf{f})$. This yields a randomized poly-time algorithm 
%
by \expref{Lemma}{Polynomial Identity}. %
For other fields, the Jacobian criterion fails due to inseparability and
%
$\AD(\F)$    %
is open.

\smallskip\noindent
\textbf{$\AM$ protocol.}
The %
Arthur-Merlin class $\AM$ is a randomized version of the class
$\NP$~\cite{AB09}.
Arthur-Merlin protocols, introduced by Babai~\cite{babai1985trading}, can be considered as a special type of interactive proof system in which the randomized poly-time verifier (Arthur) and the all-powerful prover (Merlin) have only constantly many rounds of exchange. 
$\AM$ contains interesting problems like
%
verifying that   %
two graphs are non-isomorphic.
%
$\AMcapcoAM$   %
is the class of decision problems for which both YES and NO answers can be verified by an $\AM$ protocol. It can be thought of as the randomized version of
%
$\NPcapcoNP$.    %
See~\cite{kayal2006complexity} for a few natural algebraic problems in
%
$\AMcapcoAM$.
If such a problem is $\NP$-hard (even under random reductions) then
the %
polynomial hierarchy collapses to the
%
second level,   %
\ie, 
%
$\PH=\Sigma_2$. %

In this
%
article, $\AM$ protocols  %
will only be used to distinguish whether a set $S$ is ``small'' 
or ``large.'' Formally, we refer to the Goldwasser-Sipser~\cite{goldwasser1986private} Set Lowerbound 
%
method. %
\begin{lemma}[Goldwasser-Sipser~\cite{goldwasser1986private}] \label{lem_amprotocol}
Let $m\in\N$ be given in binary. Suppose $S$ is a set whose membership can be tested in nondeterministic polynomial time and its size is promised to be either $\le m$ or $\ge 2m$. Then, the problem of deciding whether 
%
$|S| \ge 2m$ %
is in $\AM$.
\end{lemma}
 See, \eg, \cite[Chapter~8]{AB09} for a proof.
We also use the
%
Polynomial Identity Lemma  %
several times.
 The version we use is stated below.
 %
 \begin{lemma}[{Polynomial Identity Lemma~\cite{Sch80,zippel1979probabilistic,demillo1978probabilistic}}] \label{Polynomial Identity}
Let $f(x_1,\ldots,x_n)$ be a nonzero polynomial of total degree $d$ over a field $\F$. Let $S$ be a finite subset of $\F$. Then the total number of roots of $f(x_1,\ldots,x_n)$ in $S^n$ is bounded by $d|S|^{n-1}$.  
\end{lemma}

For a proof, see~\cite[Lemma~A.36]{AB09}.
 %

\smallskip\noindent
\textbf{Geometry.} Our proof requires some basic facts and definitions from classical algebraic geometry, which are listed below. One can also refer to a standard text, \eg, \cite{Har92, hartshorne2013algebraic}.  

\smallskip
Let $\aff:=\ol{\F}$ be the algebraic closure of a field $\F$. For $d\in\N^+$, write $\aff^d$ for the \emph{$d$-dimensional affine space} over $\aff$. It is defined to be the set $\aff^d$, equipped with the \emph{Zariski topology}, defined as %
\begin{definition}[Zariski topology]   %
A subset $S$ of $\aff^d$ is \emph{closed}
%
if   %
it is the set of common zeros of some
%
set   %
of polynomials in $\aff[X_1,\dots,X_d]$. For
%
a set $S$ of polynomials,
%
%
the \emph{closure} $\ol{S}$
%
is defined as  %
the smallest closed set containing $S$.
%
A set  %
$S$ is \emph{dense} in $\aff^d$  if $\ol{S}=\aff^d$.
%
%
The complement of a closed set is called an \emph{open} set.  %
\end{definition}

%
%
%

A closed set is called a \emph{hypersurface} 
if it is definable by a single polynomial.
A hypersurface is a \emph{hyperplane} if its
defining polynomial is linear.

Define $\aff^\times:= \aff\setminus\{0\}$.
Write $\proj^d$ for the $d$-dimensional projective space over $\aff$, defined to be the quotient set $(\aff^{d+1}\setminus\{(0,\dots,0)\})/\sim$
%
where   %
$(x_0,\dots,x_d)\sim (y_0,\dots,y_d)$ iff there exists $c\in\aff^\times$ such that $y_i=c x_i$ for $0\leq i\leq d$.
%
We use $(d+1)$-tuples $(x_0,\dots,x_d)$ to represent points in $\proj^d$.
%
%
Such a $(d+1)$-tuple is called a list of \emph{homogeneous coordinates}
of the point it represents.   %

The set $\proj^d$ is again equipped with the \emph{Zariski topology}.
%
\begin{definition}[Zariski topology on a projective space] %
A subset $S$ of $\proj^d$ is closed
%
if  %
it is the set of common zeros of some
%
set   %
of \emph{homogeneous} polynomials in  $\aff[X_0,\dots,X_d]$. 
\end{definition}

%
%
\begin{definition}[Variety]  %
Closed subsets of $\aff^d$ or $\proj^d$ are also called
\emph{algebraic sets} or \emph{zero sets}. 
An algebraic set is \emph{irreducible} if it cannot be written
as the union of finitely many proper algebraic sets.
An irreducible algebraic subset of an affine
%
(or projective)    %
space is also called an \emph{affine variety}
%
(\emph{projective variety}, resp.).
\end{definition}
(In some references, varieties are not required to be irreducible, but
in this  %
paper   %
we always assume
%
they are.)  %
An algebraic set $V$ can be uniquely represented as the union of finitely many varieties, and these varieties are called the \emph{irreducible components} of $V$.

%
%
Affine zero sets are in 1-1 correspondence with \emph{radical} ideals;
varieties correspond to \emph{prime} ideals.

Irreducible decomposition of an affine variety mirrors the factoring of an ideal into primary ideals.  Finally, note that the affine points are in 1-1 correspondence with \emph{maximal} ideals; this is a simple reformulation of Hilbert's Nullstellensatz.


\smallskip
The affine space $\aff^d$ may be regarded as a subset of $\proj^d$ via the map $(x_1,\dots,x_d)\mapsto (1,x_1,\dots,x_d)$. 
Then the subspace topology of $\aff^d$ induced from the Zariski topology of $\proj^d$ is just the Zariski topology of $\aff^d$.
The set $\proj^d\setminus \aff^d$ is the projective subspace of $\proj^d$ defined by $X_0=0$, called the \emph{hyperplane at infinity}.
 
For an algebraic subset $V$ of $\aff^d\subseteq\proj^d$, the smallest algebraic subset $V'$ of $\proj^d$ containing $V$ (\ie, the intersection of all algebraic subsets containing $V$) is the \emph{projective closure} of $V$, and we have $V'\cap \aff^d=V$. To see this, note that for $P=(x_1,\dots,x_d)\in \aff^d\setminus V$, there exists a polynomial $Q\in\aff[X_1,\dots,X_d]$ of degree $D\in\N$ not vanishing on $P$ (but vanishing on $V$). Then its homogenization $Q'\in\aff[X_0,\dots,X_d]$, defined by replacing each monomial
\[
  M=\prod_{i=1}^d X_i^{d_i}\qquad\text{by}\qquad  X_0^{D-\deg(M)}\prod_{i=1}^d X_i^{d_i},
\]
does not vanish on $(1,x_1,\dots,x_d)$. So, $(1,\mathbf x)\notin V'$.

\smallskip
For distinct points $P=(x_0,\dots,x_d), Q=(y_0,\dots,y_d)\in\proj^d$, write $\overline{PQ}$ for the \emph{projective line} passing through them, \ie, $\overline{PQ}$ consists of the points $(ux_0+vy_0,\dots, ux_d+vy_d)$, where $(u,v)\in\aff^2\setminus\{(0,0)\}$. 

%
\begin{definition}[Dimension and degree]  %
The \emph{dimension} of a variety $V$ is defined to be the largest integer $m$ such that there exists a chain of varieties $\emptyset\subsetneq V_0\subsetneq V_1\subsetneq\cdots\subsetneq V_m=V$. More generally, the dimension of an algebraic set $V$, denoted by $\dim V$, is the maximal dimension of its irreducible components. \Eg, we have $\dim\aff^d=\dim\proj^d=d$. 
The dimension of  the  empty set is $-1$ by convention.
%
One-dimensional    %
varieties are called \emph{curves}.

%
The \emph{degree} of a variety $V$ in $\aff^d$
%
(or in~$\proj^d$)   %
is the number of intersections of $V$ with a
%
general\footnote{We do not rigorously define the word \emph{general} here, but just remark that it essentially means ``random'' in algebraic geometry. In fact, as long as we work over a large enough field, replacing ``general'' by ``random'' would not cause any problem.} affine subspace
%
(projective subspace, resp.)  %
of dimension $d-\dim V$. 
More generally, we define the degree of an algebraic set $V$, denoted by $\deg(V)$,  to be the sum of the degrees of its irreducible components.
The degree of an algebraic subset of $\aff^d$ coincides with the degree of its projective closure in $\proj^d$.
\end{definition}

%
Suppose $V\subseteq\aff^d$ is an algebraic set, defined by polynomials $f_1,\dots,f_k$. Let $(a_1,\dots,a_d)\in\aff^d$. Then the set $\{(x_1+a_1,\dots,x_d+a_d): (x_1,\dots,x_d)\in V\}$ is called a \emph{translate} of $V$. It is also an algebraic set, defined by $f_i(X_1-a_1,\dots,X_d-a_d)$, $i=1,\dots,k$.

%
\begin{definition}[Morphism] %
Let $V\subseteq \aff^n$, $W\subseteq\aff^m$ be affine varieties. A \emph{morphism} from $V$ to $W$ is a function $f:V\to W$ that is a restriction of a polynomial map $\aff^n\to\aff^m$.
\end{definition}

A morphism $f: V\to W$ is called \emph{dominant} if $\overline{\im(f)}=W$.
The preimage of a closed subset under a morphism is closed (\ie, morphisms are \emph{continuous} in the Zariski topology).

For a polynomial map $f:\aff^n\to\aff^m$ and an affine variety $V\subseteq \aff^n$, $W:=\overline{f(V)}$ is also an affine variety (\ie, it is irreducible). To see this, assume to the contrary that $W$ is the union of two proper closed subsets $W_1$ and $W_2$. By the definition of closure, $f(V)$ is not contained in either $W_1$ or $W_2$, \ie, it intersects both. Then $f^{-1}(W_1)\cap V$ and $f^{-1}(W_2)\cap V$ are two proper closed subsets of $V$, and their union is $V$. This contradicts the irreducibility of $V$.

%
The \emph{graph} $\Gamma_f$ of a morphism $f$ is the set $\{(x,f(x)): x\in V\}\subseteq V\times W\subseteq \aff^n\times \aff^m$.
Here $V\times W=\{(x,y): x\in V,y\in W\}$ denotes the \emph{product} of $V$ and $W$, which is a subvariety of the $(n+m)$-dimensional affine space $\aff^n\times \aff^m\cong \aff^{n+m}$.
Note the graph $\Gamma_f$ is closed in $\aff^n\times \aff^m$: 
Suppose $f$ sends $x\in V$ to $(f_1(x),\dots,f_m(x))\in\aff^m$, where $f_i\in\aff[X_1,\dots,X_n]$ for $i\in [m]$.
And suppose $V$ is defined by an ideal $I\subseteq \aff[X_1,\dots,X_n]$. Then $\Gamma_f$ is defined by the ideal of $\aff[X_1,\dots,X_n,Y_1,\dots,Y_m]$ generated by $I$ and the polynomials $Y_i-f_i(X_1,\dots,X_n)$, $i=1,\dots,m$. 

%

\section{Algebraic dependence testing: Proof of \expref{Theorem}{thm_amcoam} }\label{sec-dep}
%

Given $f_1,\dots,f_m\in \F_q[x_1,\dots,x_n]$, we want to decide if they are algebraically dependent. For this
%
problem, $\AD(\F_q)$,    %
we could assume, with some preprocessing, that $m=n$. For, $m>n$ means that
%
it is  %
a YES instance. If $m < n$ then we could apply a ``random'' linear map on the variables to reduce 
%
$n$ %
to $m$, preserving the YES/NO instances. Also, the transcendence degree does not change when we move to the algebraic closure $\ol{\F}_q$. The details can be found in~\cite[Lemmas~2.7--2.9]{pandey2016algebraic}. So, we assume the input instance to be $\mathbf f := \inbrace{f_1,\dots,f_n}$ with nonconstant polynomials.
 
 In the following, let $D:=\prod_{i\in [n]}\deg(f_i) \,>0$ and $D':=\max_{i\in [n]} \deg(f_i) \,>0$. Let $d\in\N^+$ and $q'=q^d$. The value of $d$ will be determined later.
 Let $f:\F_{q'}^n\to\F_{q'}^n$ be the polynomial map $a\mapsto (f_1(a),\dots,f_n(a))$. For $b=(b_1,\dots,b_n)\in \F_{q'}^n$, denote by $N_b$ the size of the preimage
 \[
   f^{-1}(b)=\left\{\mathbf x \in \F_{q'}^n \,\bigm|\, f(\mathbf x) = b \right\}.
 \]

Define $\aff:=\ol{\F}_q$ and $\overline{N}_b:= \#\{\mathbf x \in \aff^n \,\mid\, f_i(\mathbf x) = b_i, \, \text{ for all } i\in[n] \}$ which might be $\infty$.
Let $Q\in\F_q[y_1,\dots,y_n]$ be a nonzero annihilator, of minimal degree, of $f_1,\dots,f_n$. If it exists then $\deg(Q)\leq D$ by Perron's bound. 

\subsection{\texorpdfstring{$\AM$}{AM} protocol}

First, we study the independent case.

%
%
\begin{lemma}[Dimension-zero preimage]\label{lem_finite}
Suppose $\mathbf f$
%
is  %
independent. Then $\overline{N}_{f(a)}$ is \emph{finite} for all but at most $(nDD'/q')$-fraction of $a\in\F_{q'}^n$.
\end{lemma}
\begin{proof}
For $i\in [n]$, let $G_i\in\F_q[z, y_1,\dots,y_n]$ be the annihilator of $\inbrace{x_i, f_1,\dots,f_n}$. We have $\deg(G_i)\leq D$ by Perron's bound. 
Consider $a\in\F_{q'}^n$ such that  $G'_i(z):= G_i(z, f_1(a),\dots,f_n(a))\in\F_q[z]$ is a nonzero polynomial for every $i\in [n]$.
We claim that $\overline{N}_{f(a)}$ is finite for such $a$. 

To see this, note that for any $b=(b_1,\dots,b_n)\in \aff^n$ satisfying the equations $f_i(b)=f_i(a)$, $i\in[n]$, we have
\[
0=\, G_i(b_i, f_1(b), \dots, f_n(b)) \,=\, G_i(b_i, f_1(a),\dots, f_n(a)) \,=\, G'_i(b_i), \quad\forall i\in [n] \,.
\]
Hence, each $b_i$ is a root of $G'_i$. It follows that $\overline{N}_{f(a)} \leq \prod_{i\in[n]} \deg(G'_i) < \infty$, as claimed. 

It remains to prove  that the number of $a\in\F_{q'}^n$ satisfying $G'_i=0$, for some index $i\in [n]$, is bounded by $nDD'q'^{-1}\cdot q'^{n}$.
Fix $i\in [n]$. Suppose $G_i=\sum_{j=0}^{d_i} G_{i,j} z^j$, where $d_i:=\deg_z(G_i)$ and $G_{i,j}\in \F_q[y_1,\dots,y_n]$, for $0\leq j\leq d_i$.
The leading coefficient $G_{i,d_i}$ is a nonzero polynomial. As $f_1,\dots,f_n$ are algebraically independent, the polynomial $G_{i,d_i}(f_1,\dots,f_n)\in\F_q[x_1,\dots,x_n]$ is also nonzero.
Its degree is $\le D'\deg(G_{i,d_i}) \leq D'\deg(G_i) \leq DD'$.
By 
%
\expref{Lemma}{Polynomial Identity},  %
for all but at most $(DD'/q')$-fraction of $a\in \F_{q'}^n$, we have 
%
\[
G_{i,d_i}(f_1(a),\dots,f_n(a))\neq 0,
\]  %
which   implies  
\[
G'_i(z) \,=\, G_i(z, f_1(a),\dots,f_n(a)) \,=\, \sum_{j=0}^{d_i} G_{i,j}(f_1(a),\dots,f_n(a)) z^j \,\neq 0 \,.
\]
The claim now follows from the union bound.
\end{proof}

We need the following affine version of B\'ezout's Theorem. Its proof can be found in~\cite[Theorem~3.1]{Sch95}.
\begin{theorem}[B\'ezout's]
 Let $g_1,\dots,g_n\in \aff[x_1,\dots,x_n]$. Then the number of common zeros of $g_1,\dots,g_n$ in $\aff^n$ is either infinite, or at most $\prod_{i\in[n]} \deg(g_i)$.
\end{theorem}

Combining \expref{Lemma}{lem_finite} with B\'ezout's Theorem, we obtain
\begin{lemma}[Small preimage]\label{lem_bpre}
Suppose $\mathbf f$
%
is  %
independent.
Then $N_{f(a)}\leq D$ for all but at most $(nDD'/q')$-fraction of $a\in\F_{q'}^n$.
\end{lemma}

Next, we study the dependent case (with an annihilator $Q$).

\begin{lemma}[Large preimage]\label{lem_coamdep}
Suppose $\mathbf f$
%
is  %
dependent.
Then for $k>0$, we have $N_{f(a)}>k$ for all but at most $(kD/q')$-fraction of $a\in \F_{q'}^n$.
\end{lemma}
\begin{proof}
 Let $\im(f):= f(\F_{q'}^n)$ be the image of the map. Note that $Q$ vanishes on all the points in $\im(f)$. So, $|\im(f)| \leq Dq'^{n-1}$ by 
 %
 \expref{Lemma}{Polynomial Identity}. %
 
Let $B:= \{b\in \im(f): N_b\leq k\}$ be the ``bad'' images. We can estimate the bad domain points as, 
\[
\#\{a\in \F_{q'}^n: N_{f(a)}\leq k\} \,=\, \#\{a\in \F_{q'}^n: f(a)\in B\} \,\le\,  k|B| \le k|\im(f)| \,\le\, k D q'^{n-1} \,.
\]
which proves the lemma.
\end{proof}


\begin{theorem}[$\AM$]\label{thm_am}
Testing algebraic dependence of $\mathbf f$ is in $\AM$.
\end{theorem}
\begin{proof}
Fix $q'=q^d>4nDD'+4kD$ and $k:=2D$. Note that $d$ will be polynomial in the input size.
For an $a\in \F_{q'}^n$, consider the set 
$f^{-1}(f(a)) :=$ $\{\mathbf x \in \F_{q'}^n \,\mid\, f(\mathbf x) = f(a) \}$.

%
By \expref{Lemma}{lem_bpre} and \expref{Lemma}{lem_coamdep}, when %
Arthur picks $a$ randomly, with high probability, $|f^{-1}(f(a))|= N_{f(a)}$ is more than $2D$ in the dependent case while $\le D$ in the independent case. Note that an upper bound on $\prod_{i\in [n]}\deg(f_i)$ can be deduced from the size of the input circuits for 
%
the polynomials $f_i$; %
thus, we know $D$. Moreover, containment in $f^{-1}(f(a))$ can be tested in 
%
$\Pclass$. %
Thus, by \expref{Lemma}{lem_amprotocol},
%
$\AD(\F_q)$    %
is in $\AM$.
\end{proof}


\subsection{\texorpdfstring{$\coAM$}{coAM} protocol}
We first study the independent case.
%

\begin{lemma}[Large image]\label{lem-coam-1}
Suppose $\mathbf f$
%
is  %
independent.
Then $N_b>0$ for at least $(D^{-1} - nD'q'^{-1})$-fraction of $b\in \F_{q'}^n$.
\end{lemma}
\begin{proof}
Let $S:= \{a\in \F_{q'}^n: N_{f(a)}\leq D\}$.
Then $|S|\geq (1-n DD'q'^{-1})\cdot q'^{n}$ by \expref{Lemma}{lem_bpre}.
As every $b\in f(S)$ has at most $D$ preimages in $S$ under $f$, we have 
$|f(S)| \,\ge\, |S|/D \,\ge\, (D^{-1}-nD'q'^{-1})\cdot q'^{n} $.
This proves the lemma since $N_b>0$ for all $b\in f(S)$.
\end{proof}

Next, we study the dependent case.

\begin{lemma}[Small image]\label{lem-coam-2}
Suppose $\mathbf f$
%
is  %
dependent.
Then $N_b=0$ for all but at most $(D/q')$-fraction of $b\in \F_{q'}^n$.
\end{lemma}
\begin{proof}
%
By definition, %
$N_b>0$ iff $b\in \im(f):=f(\F_{q'}^n)$. It was shown in the proof of \expref{Lemma}{lem_coamdep} that $|\im(f)| \le Dq'^{n-1}$. The lemma follows.
\end{proof}

\begin{theorem}[$\coAM$]\label{thm_coam}
Testing algebraic dependence of $\mathbf f$ is in $\coAM$.
\end{theorem}
\begin{proof}
Fix $q'=q^d \,> D(2D+nD')$. Note that $d$ will be polynomial in the input size.
For $b\in \F_{q'}^n$, consider the set 
$f^{-1}(b) :=$ $\{\mathbf x \in \F_{q'}^n \,\mid\, f(\mathbf x) = b \}$ of size $N_b$.

Define $S:= \im(f)$. Note that %
$b\in \F_{q'}^n$
%
satisfies   %
$N_b>0$ iff $b\in S$.
%
%
%
Thus, by \expref{Lemma}{lem-coam-1},
$|S|\ge (D^{-1} - nD'q'^{-1})q'^n\,> 2Dq'^{n-1}$ 
when $\mathbf f$ is independent, and
by \expref{Lemma}{lem-coam-2}, $|S|\le Dq'^{n-1}$
when $\mathbf f$ is dependent.

Note that an upper bound on $\prod_{i\in [n]}\deg(f_i)$ can be deduced from the size of the input circuits for 
%
the polynomials $f_i$; %
thus, we know $Dq'^{n-1}$. Moreover, containment in $S$ can be tested in $\NP$. Thus, by \expref{Lemma}{lem_amprotocol},
%
$\AD(\F_q)$    %
is in $\coAM$.
\end{proof}

\begin{proof}[Proof of \expref{Theorem}{thm_amcoam}]
The statement immediately follows from \expref{Theorem}{thm_am} and \expref{Theorem}{thm_coam}.
\end{proof}


%

\section{Approximate polynomials satisfiability: Proof of \expref{Theorem}{thm-aps}} \label{sec:aps}
%

\expref{Theorem}{thm-aps} is proved in two parts. First, we show that $\APS$ is equivalent to $\AnnAtZero$ 
%
problem, %
which means that it is $\NP$-hard~\cite{Kay09}. Next, we utilize the  beautiful underlying geometry to devise a $\PSPACE$ algorithm.


\subsection{\texorpdfstring{$\APS$}{APS} is equivalent to \texorpdfstring{$\AnnAtZero$}{AnnAtZero}}

The proof that $\APS$ is equivalent to $\AnnAtZero$ is classical and follows from two statements in algebraic geometry and commutative algebra, which we state as \expref{Lemma}{lem_geom} and \expref{Theorem}{thm_approx}.
The first statement (\expref{Lemma}{lem_geom}) gives a geometric reinterpretation of $\AnnAtZero$. This can be seen as a direct application of the ideal-variety correspondence. 
The second statement (\expref{Theorem}{thm_approx})  is also classical. It was essentially proved in~\cite{LL89} and was also discussed in~\cite[Lemma~20.28]{burgisser2013algebraic} and~\cite[Page~37:53]{grochow2018circuit}.

%

Let $\aff$ be the algebraic closure of $\F$. Note that for the given polynomials $\mathbf f:= \inbrace{f_1,\dots,f_m}$ in $\F[\mathbf x]$, there is an annihilator over $\F$ with 
%
a nonzero %
constant term iff there is an annihilator over $\aff$ with 
%
a nonzero %
constant term. This is because if $Q$ is an annihilator over $\aff$ with 
%
a nonzero %
constant term, w.l.o.g.\ $1$, then by basic linear algebra, the linear system defined by the equation $Q(\mathbf f)=0$, in terms of the unknown coefficients of $Q$, would also have a solution in $\F$. Thus, there is an annihilator over $\F$ with constant term $1$. This proves that it suffices to solve $\AnnAtZero$ over the algebraically closed field $\aff$. This provides us with a better geometry. 

Write $f:\aff^n\to \aff^m$ for the polynomial map sending a point
\[
  x=(x_1,\dots,x_n)\in \aff^n\quad\text{to}\quad (f_1(x),\dots,f_m(x))\in \aff^m\,.
  \]
  For a subset $S$ of an affine or projective space, write
  $\overline{S}$ for its Zariski closure in that space. We will use
  $O$ to denote the origin $\mathbf 0$ of an affine space.

%
 
The following lemma reinterprets $\AnnAtZero$ in a geometric way.
 \begin{lemma}[classical]\label{lem_geom}
 The constant term of every annihilator for $\mathbf f$ is zero iff $O\in \overline{\im(f)}$. %
 \end{lemma}
 \begin{proof}
 %
 Note that %
 $Q\in \aff[Y_1,\dots,Y_m]$ vanishes on $\im(f)$ iff $Q(\mathbf f)$ vanishes on $\aff^n$, which holds iff $Q(\mathbf f)=0$, \ie, $Q$ is an annihilator for $\mathbf f$. So $\overline{\im(f)}=V(I)$, where the ideal $I\subseteq \aff[Y_1,\dots,Y_m]$ consists of the annihilators for $\mathbf f$. Also note that $\{O\}=V(\m)$, where $\m$ is the maximal ideal $\langle Y_1,\dots, Y_m\rangle$.

Let us study the condition $O\in \overline{\im(f)}$.
 By the ideal-variety correspondence, $\{O\} \,=\, V(\m) \,\subseteq\, \overline{\im(f)} \,=\, V(I)$ is equivalent to $I\subseteq \m$, \ie,  $Q\bmod \m=0$ for  $Q\in I$. But $Q\bmod \m$ is just the constant term of the annihilator $Q$. Hence, we have the equivalence.
 \end{proof}
 
%
%
As a special case,
%
the above lemma proves that
whenever $\mathbf f$
%
is  %
algebraically \emph{independent}, we have $\aff^m=\ol{\im(f)}$. \Eg, $f_1=X_1$ and $f_2=X_1 X_2-1$. 
%

In general, %
$\im(f)$ is not necessarily closed in the Zariski topology. 
\begin{exmp}\label{exmp_closure}
Let $n=2$, $m=3$. Consider $f_1=f_2=X_1$ and $f_3=X_1 X_2-1$. The annihilators are multiples of $(Y_1-Y_2)$, which means by \expref{Lemma}{lem_geom} that $O\in \ol{\im(f)}$. But there is no solution to $f_1=f_2=f_3=0$, \ie, $O\notin \im(f)$.
\end{exmp}

\noindent
\textbf{Approximation.} Although $O\in \overline{\im(f)}$ is not equivalent to the existence of a solution $x\in \aff^n$ to $f_i=0$, $i\in[m]$, it is equivalent to the existence of an  ``approximate solution'' $x\in \aff[\epsilon,\epsilon^{-1}]^n$, which is a tuple of Laurent polynomials in a formal variable $\epsilon$. The formal statement is as follows. W.l.o.g.\ we assume $\mathbf f$ to be $m$ nonconstant polynomials.

\begin{theorem}[classical]\label{thm_approx}
$O\in \overline{\im(f)}$ iff there exists $x=(x_1,\dots,x_n)\in \aff(\epsilon)^n$ such that $f_i(x)\in \epsilon \aff[\epsilon]$, for all $i\in [m]$. Moreover, when such $x$ exists, it may be chosen such that
\[
x_i \,\in\,\, \epsilon^{-\Delta} \aff[\epsilon] \,\cap\, \epsilon^{\Delta'} \aff[\epsilon^{-1}] \,=\, \left\{\sum_{j=-\Delta}^{\Delta'} c_j \epsilon^j: c_j\in \aff\right\}, \quad i\in[n],
\]
where $\Delta:=\prod_{i\in[m]} \deg(f_i) \,>0$ and $\Delta':=(\max_{i\in[m]} \deg(f_i))\cdot \Delta \,>0$.
\end{theorem}

For example, choosing $x=(x_1,x_2):=(\epsilon, 1/\epsilon)$ in \expref{Example}{exmp_closure},  we have $f_1(x)=f_2(x)=\epsilon$ and $f_3(x)=0$. Thus \expref{Theorem}{thm_approx} gives another way of seeing  $O\in \ol{\im(f)}$ in \expref{Example}{exmp_closure}.

As mentioned previously, \expref{Theorem}{thm_approx} was  essentially proved  in~\cite{LL89}. We include a proof here for the sake of completeness.

First, we recall a tool to reduce the domain from a variety to a curve, proven in~\cite{LL89}.

\begin{lemma}
%
[{\cite[Proposition~1]{LL89}}] %
\label{lem_aux1}
Let $V\subseteq \aff^n$, $W\subseteq \aff^m$ be affine varieties, $\varphi: V\to W$ dominant, and $t\in W\setminus \varphi(V)$.
Then there exists a curve $C\subseteq \aff^n$ such that $t\in\overline{\varphi(C)}$ and $\deg(C)\leq \deg(\Gamma_\varphi)$, where $\Gamma_\varphi$ denotes the graph of $\varphi$ embedded in $\aff^n\times \aff^m$.
\end{lemma}

Next,~\cite{LL89} essentially shows that in the case of a curve one can approximate the preimage of $f$ by using a \emph{single} formal variable $\epsilon$ and working in $\aff(\epsilon)$. Here $\aff[[\epsilon]]$ denotes a formal power series in $\epsilon$ with coefficients from $\aff$.

\begin{lemma}
%
[{\cite[Corollary of Proposition~3]{LL89}}] %
\label{lem_aux2}
Let $C\subseteq \aff^n$ be an affine curve.  Let $f: C\to \aff^m$ be a morphism sending $x\in C$ to $(f_1(x),\dots,f_m(x))\in \aff^m$, where $f_1,\dots,f_m\in \aff[X_1,\dots,X_n]$. Let $t=(t_1,\dots,t_m)\in \overline{f(C)}$. Then there exists  $p_1,\dots,p_n\in \epsilon^{-\deg(C)}\aff[[\epsilon]]$ such that $f_i(p_1,\dots,p_n)-t_i \,\in\, \epsilon \aff[[\epsilon]]$ , for all $i\in [m]$.
\end{lemma}

Finally, we use the above two lemmas to prove the connection of $\APS$ with $O\in \overline{\im(f)}$, and hence with $\AnnAtZero$ (by \expref{Lemma}{lem_geom}).

\begin{proof}[Proof of \expref{Theorem}{thm_approx}]
First assume that an $x$, satisfying the conditions in \expref{Theorem}{thm_approx}, exists. Pick such an $x$. If $\mathbf f$
%
is  %
algebraically independent then by \expref{Lemma}{lem_geom} we have that $\aff^m=\overline{\im(f)}$ and we are done. So, assume that there is a nonzero annihilator $Q$ for $\mathbf f$.
We have $Q(f_1(x),\dots,f_m(x))=0\in\epsilon \aff[\epsilon]$. On the other hand, as $f_i(x)\in \epsilon \aff[\epsilon]$, for all $i\in [m]$; we deduce that $Q(f_1(x),\dots,f_m(x))\bmod \epsilon \aff[\epsilon]$ is $Q(\mathbf 0)$, which is the constant term of $Q$. So it equals zero. By \expref{Lemma}{lem_geom}, we have $O\in \overline{\im(f)}$ and again we are done.

Conversely, assume $O\in \overline{\im(f)}$ and we will prove that $x$ exists.
If $O\in \im(f)$, then we can choose $x\in \aff^n$ and we are done. So assume $O\in  \overline{\im(f)}\setminus \im(f)$.
Regard $f$ as a dominant morphism from $\aff^n$ to $\overline{\im(f)}$. Its graph $\Gamma_f$ is cut out in $\aff^n\times \aff^m$ by $Y_i-f_i(X_1,\dots,X_n)$, $i\in[m]$. So $\deg(\Gamma_f)\leq \prod_{i=1}^m \deg(f_i) = \Delta$ by B\'ezout's Theorem.

By \expref{Lemma}{lem_aux1}, there exists a curve $C\subseteq \aff^n$ such that $O\in \overline{f(C)}$ and $\deg(C)\leq \deg(\Gamma_f)\leq \Delta$. Pick such a curve $C$. Apply \expref{Lemma}{lem_aux2} to $C$, $f|_C$ and $O$, and let $p_1,\dots,p_n\in \epsilon^{-\deg(C)} \aff[[\epsilon]]\subseteq \epsilon^{-\Delta} \aff[[\epsilon]]$ be as given by the lemma.
Then $f_i(p_1,\dots,p_n)\in \epsilon \aff[[\epsilon]]$, for all $i\in [m]$. 

For $i\in [n]$, let $x_i$ be the Laurent polynomial obtained from $p_i$ by truncating the terms of degree greater than $\Delta'$.
When evaluating $f_1,\dots,f_m$, at $(p_1,\dots,p_n)$, such truncation does not affect the coefficient of $\epsilon^k$ for $k\leq 0$ by the choice of $\Delta'$.
So $f_i(x_1,\dots,x_n)\in \epsilon \aff[\epsilon]$,  for all $i\in [m]$.
\end{proof}

\noindent

%
\begin{remark}   %
The lower bound $-\Delta=-\prod_{i=1}^m \deg(f_i)$ for the least degree of $x_i$ in $\epsilon$ can be achieved up to a  factor of $1+o(1)$. Consider the polynomials $f_1=f_2=X_1$, $f_3=X_1^{d-1}X_2-1$, and $f_i=X_{i-2}^d-X_{i-1}$ for $i=4,\dots,m$, where $m=n+1$. Then we are forced to choose $x_1\in  \epsilon \aff[\epsilon]$ and $x_i \,\in\, \epsilon^{-(d-1)d^{i-2}}\cdot\aff[\epsilon^{-1}] $, for $i=2,\dots, n$. So the least degree of $x_n$ in $\epsilon$  is at most $-(d-1)d^{n-2}$, while $-\Delta=-d^{n-1}$.
\end{remark}

\subsection{Putting \texorpdfstring{$\APS$}{APS} in \texorpdfstring{$\PSPACE$}{PSPACE}}
%

Owing to the exponential upper bound on the precision (= degree w.r.t.\ $\epsilon$) shown in \expref{Theorem}{thm_approx}, one expects to solve $\APS$ in
%
$\EXPSPACE$   %
only. Surprisingly, in this section, we give a $\PSPACE$ algorithm. This we do by reducing the general $\AnnAtZero$ instance to a very special instance, that is easy to solve.

Let $\aff$ be the algebraic closure of the field $\F$. Let $f_1,\dots,f_m\in \F[X_1,\dots,X_n]$ be given. 
Denote by $k$ the transcendence degree of $\F(f_1,\dots,f_m)/\F$. Computing $k$ can be done in $\PSPACE$ using linear algebra~\cite{Plo05,C76,BGH82, Mul87}. 
We assume $k<m-1$, since the cases $k=m-1$ and $k=m$ are again easy. In the case $k=m$, the input instances are always in $\APS$ since $\overline{\im(f)}=\aff^m$. And in the case $k=m-1$, the ideal of the annihilators is a principal ideal, and hence has a unique generator (up to scaling). The degree of this generator is at most $\prod_{i=1}^m \deg(f_i)$. Thus checking whether it has a nonzero constant term can be solved in $\PSPACE$ by solving an exponential sized linear system of equations using~\cite{C76, BGH82, Mul87}. 

We reduce the number of polynomials from $m$ to $k+1$ as follows:
Fix a finite 
%
set %
$S\subseteq \F$,  and choose $c_{i,j}\in S$ at random for $i\in [k+1]$ and $j\in [m]$. For this to work, we need a large enough $S$ and $\F$.
For $i\in [k+1]$, let $g_i:= \sum_{j=1}^m c_{i,j} f_j$. 

Let  $\delta:= (k+1)(\max_{i\in [m]}\deg(f_i))^k/|S|$. Our algorithm is immediate once we prove the following claim.
%
\begin{theorem}[Random reduction]\label{thm-randReduct}
It holds, with probability $\ge (1-\delta)$, that  

\smallskip\noindent
(1) the transcendence degree of $\F(g_1,\dots,g_{k+1})/\F$ equals $k$, and

\noindent
(2) the constant term of every annihilator for $g_1,\dots,g_{k+1}$ is zero  iff the constant term of  every annihilator for $f_1,\dots,f_{m}$ is zero.
\end{theorem}

First, we reformulate the two items of \expref{Theorem}{thm-randReduct} in a geometric way, and later we will analyze the error probability.

For $d\in\N$, denote by $\aff^d$
%
(and by~$\proj^d$)
the $d$-dimensional affine space
%
(projective space, resp.) over $\aff:=\ol{\F}$.
Let $f: \aff^n\to \aff^m$ (resp.~$g: \aff^n\to \aff^{k+1}$)
be the polynomial map sending $x$ to $(f_1(x),\dots,f_m(x))$
(resp.~$(g_1(x),\dots,g_{k+1}(x))$). Let $O$ and $O'$ be the
origin of $\aff^m$ and that of $\aff^{k+1}$ respectively.
Define the affine varieties $V:= \overline{\im(f)}\subseteq \aff^m$
and $V':= \overline{\im(g)}\subseteq \aff^{k+1}$. 
Then $\dim V= \trdeg(\mathbf f) = k$.

Let $\pi: \aff^m\to \aff^{k+1}$ be the linear map sending $(x_1,\dots,x_m)$ to $(y_1,\dots,y_{k+1})$ where $y_i= \sum_{j=1}^m c_{i,j} x_j$.
Then $g= \pi\circ f$ and $V'= \overline{\pi(V)}$.\footnote{To see $V'\supseteq\overline{\pi(V)}$, note that $\pi^{-1}(V')$ contains $\im(f)$ and is closed, and hence contains $V=\overline{\im(f)}$.}
 Now (1) of  \expref{Theorem}{thm-randReduct} is equivalent to $\dim V'=k$, and (2) is equivalent to $O'\in V'$ iff  $O\in V$.
\[
\begin{tikzcd}[column sep=large]
\aff^n \arrow{r}{f} \arrow[swap]{rd}{g}
& V= \overline{\im(f)}  \arrow{r}{\subseteq} \arrow{d}{\pi|_V} & \aff^m \arrow{d}{\pi}\\
& V'= \overline{\im(g)} \arrow{r}{\subseteq} & \aff^{k+1}
\end{tikzcd}
\] 
 We will give sufficient conditions of (1) and (2) in terms of incidence properties.
 Note that $O\in V$ implies $O'\in V'$, since $\pi(O)=O'$. Now suppose $O\not\in V$. Let $W:= \pi^{-1}(O')$, which is a linear  subspace of $\aff^m$.
  Then $O'\not\in \pi(V)$ iff $V\cap W=\emptyset$. However, $V\cap W=\emptyset$ does not imply $O'\not\in V'$, as $V$ may  ``get infinitesimally close to $W$''  without actually meeting $W$, so that $O'\in\overline{\pi(V)}=V'$. This is shown by the following 
 %
   example. %
 %

  
    \begin{exmp}\label{exmp_reduction}
Let $m=4$ and $(f_1,f_2,f_3,f_4)=(X_1,X_2,X_1 X_2-1, X_1+X_2)$. Then $k:=\trdeg(\mathbf f)= 2$.
  Let $(g_1,g_2,g_3)=(f_1,f_3,f_1+f_2-f_4)=(X_1,X_1 X_2-1,0)$.
  Suppose $\aff^m$ has coordinates $Y_1,\dots,Y_4$ and $\aff^{k+1}$ has coordinates $Z_1,\dots,Z_3$.

Then $V\subseteq\aff^m$ is defined by $Y_1Y_2-Y_3-1=0$ and $Y_1+Y_2-Y_4=0$, and $W$ is defined by $Y_1=0$, $Y_3=0$, and $Y_2-Y_4=0$.
So $V\cap W=\emptyset$.  But $V'\subseteq\aff^{k+1}$ is the plane $Z_3=0$, which contains the origin.
  \end{exmp}
  
  To overcome the above problem, we consider projective geometry instead of affine geometry.
  Suppose $\aff^m$ have coordinates $X_1,\dots,X_m$ and $\proj^m$ have homogeneous coordinates $X_0,\ldots,X_m$.
  Regard $\aff^m$ as a dense open subset of $\proj^m$ via $(x_1,\dots,x_m)\mapsto (1,x_1,\dots,x_m)$. Then $H:= \proj^m\setminus \aff^m \,\cong \proj^{m-1}$ is the \emph{hyperplane at infinity}, defined by $X_0=0$.
Denote by $V_c$
%
(and $W_c$)   %
the \emph{projective closure} of $V$
%
($W$, resp.)   %
in $\proj^m$. Then $V = V_c\cap \aff^m$.
  Let $W_H:= W_c\cap H$, which is a projective subspace of $H$. 
  
  For distinct points $P,Q\in\proj^m$, write $\overline{PQ}$ for the projective line passing through them.  The following two lemmas  provide sufficient conditions for $\dim V'=k$ and $O'\not\in V'$, respectively.
%
%
\begin{lemma}\label{lem_projcriterion}  
 If $V_c\cap W_H=\emptyset$, then  $\dim V'=k$.
\end{lemma}

\begin{proof}
Assume $\dim V'<k$. Choose $P\in\pi(V)$. The dimension of $\pi^{-1}(P)\cap V$ is at least $\dim V-\dim V' \,\geq 1$~\cite[Theorem~11.12]{Har92}.
Denote by $Y$ and $Z$  the projective closure of $\pi^{-1}(P)$ and that of $\pi^{-1}(P)\cap V$ in $\proj^m$, respectively.
%
Then $Z\subseteq Y\cap V_c$. As $\dim Z=\dim \pi^{-1}(P)\cap V\geq 1$ and $\dim H=m-1$, we have $Z\cap H\neq\emptyset$~\cite[Proposition~11.4]{Har92}. 

As $\pi$ is a linear map, $\pi^{-1}(P)=Y\cap\aff^m$ is a translate of $\pi^{-1}(O')= W =W_c\cap\aff^m$. It is well known that two projective subspaces $W_1,W_2\not\subseteq H$ have the same intersection with $H$ iff $W_1\cap \aff^m$ and $W_2\cap \aff^m$ are translates of each other.\footnote{Indeed, $W_i\cap \aff^m$ is defined by linear equations $\sum_{j=1}^m a_{j,t} X_j +a_{0,t}=0$ iff $W_i\cap H$ is defined by homogeneous linear equations $X_0=0$ and $\sum_{j=1}^m a_{j,t} X_j=0$. So the constant terms $a_{0,t}$ do not matter.}
So, $Y\cap H=W_c\cap H=W_H$. Therefore,
$V_c\cap W_H \,=\, V_c\cap Y\cap H \,\supseteq\, Z\cap H \,\neq \emptyset$.
\end{proof}

\begin{lemma}\label{lem_projcriterion2}
If  $V_c\cap W_c=\emptyset$, then $O'\not\in V'$. 
\end{lemma}
\begin{proof}

 Assume to the contrary that   $V_c\cap W_c=\emptyset$ but $O'\in V'$. We will derive a contradiction. 
As $W_H\subseteq W_c$, we have  $V_c\cap W_H=\emptyset$ and hence $\dim V'=k$ by \expref{Lemma}{lem_projcriterion}.

Denote by $J(V_c, W_H)$ the \emph{join} of $V_c$ and $W_H$, which is defined to be the union of the projective lines $\overline{PQ}$, where $P\in V_c$ and $Q\in W_H$. It is known that $J(V_c, W_H)$, as the join of two \emph{disjoint} projective subvarieties, is again a projective subvariety of $\proj^m$~\cite[Example~6.17]{Har92}. 

For $P\in V$,  denote by $W_P$ the unique translate of $W$ containing $P$. We need the following two 
%
claims. %

\begin{claim}\label{claim_join}
$J(V_c, W_H)\cap \aff^m \,=\, \bigcup_{P\in V} W_P$.
\end{claim}
\begin{proof}[Proof of \expref{Claim}{claim_join}]
Consider $P\in V_c$ and $Q\in W_H$. If $P\in H$, the line $\overline{PQ}$ lies in $H$ and does not meet $\aff^m$. Now suppose $P\in V_c\setminus H=V$. Then $\overline{PQ}$ meets $\overline{OQ}$ at the point $Q$. So $\overline{PQ}\cap \aff^m$ is a translate of $\overline{OQ}\cap \aff^m \,\subseteq\, W_c\cap\aff^m = W$. It follows that $\overline{PQ}\cap \aff^m\subseteq W_P$. This proves  $J(V_c, W_H)\cap \aff^m\subseteq \bigcup_{P\in V} W_P$.

Conversely, let $P\in V$.   Let $\ell_P$ be an affine line contained in $W_P$ and passing through $P$ (note that $W_P$ is the union of such lines). Then $\ell_P$ is a translate of an affine line $\ell\subseteq W$. As $\ell_P$ and $\ell$ are translates of each other, their projective closures intersect $H$ at the same point $Q$. We have $Q\in \ell\cap H \subseteq W_H$. So $\ell_P=\overline{PQ}\cap \aff^m\subseteq J(V_c, W_H)\cap \aff^m$.
\end{proof}

%
%
%
%

\begin{claim}\label{claim_join2}
  $J(V_c,W_H)\cap \aff^m=\pi^{-1}(V')$\,.
  \end{claim} 
 \begin{proof}[Proof of \expref{Claim}{claim_join2}]
  As $\pi$ is a linear map, \expref{Claim}{claim_join}  implies $J(V_c,W_H)\cap \aff^m\subseteq\pi^{-1}(V')$. We prove the other direction by comparing dimensions.
It is known that for two \emph{disjoint} projective subvarieties $V_1$ and $V_2$,  $\dim J(V_1,V_2)=\dim V_1+\dim V_2+1$~\cite[Proposition~11.37 and Exercise~11.38]{Har92}. Therefore,
\[
\dim J(V_c, W_H)=\dim V_c +\dim W_H+1=\dim V+\dim W=k+\dim W \,.
\]
So, $\dim J(V_c,W_H) \cap \aff^m \,= k+\dim W$.
On the other hand, we have $\pi^{-1}(V')\cong V'\times W$. So $\dim \pi^{-1}(V')=\dim V'+\dim W=k+\dim W$. Now $J(V_c,W_H)\cap \aff^m$ and $\pi^{-1}(V')$ are (irreducible) affine varieties of the same dimension, and one is contained in the other. So they must be equal. This proves the claim.
\end{proof}

So  $\pi^{-1}(V')=\bigcup_{P\in V} W_P$ by \expref{Claim}{claim_join} and \expref{Claim}{claim_join2}.
As $O'\in V'$, we have $W=\pi^{-1}(O')\subseteq \pi^{-1}(V')=\bigcup_{P\in V} W_P$. So $W_P=W$ for some $P\in V$, since $W$ is a linear space. But then $P\in V\cap W_P=V\cap W\subseteq V_c\cap W_c$, contradicting the assumption $V_c\cap W_c=\emptyset$.
 \end{proof}
 %
\noindent
%
\begin{remark}   %
The converse of \expref{Lemma}{lem_projcriterion2} is false, as shown by the following example.  %

\begin{exmp}\label{exmp-projcrit}
Consider \expref{Example}{exmp_reduction} but choose $f_4$ to be $X_1+X_2+1$ instead of $X_1+X_2$. Now we have $g_3=1$, $V$ is defined by $Y_1Y_2-Y_3-1=0$ and $Y_1+Y_2-Y_4+1=0$, and $V'$ is the plane $Z_3=1$. So $O'\not\in V'$. 

On the other hand, suppose $\proj^m$ has coordinates $Y_0,\dots,Y_4$. Then $V_c\cap H$ is defined by $Y_0=Y_1Y_2=Y_1+Y_2-Y_4=0$, and $W_H$ is defined by $Y_0=Y_1=Y_2-Y_4=Y_3=0$. So $(0,0,1,0,1)\in V_c\cap W_H\subseteq V_c\cap W_c$.
\end{exmp}
\end{remark}

\smallskip\noindent
\textbf{Error probability.} It remains to bound the probability of failure of the conditions $V_c\cap W_H=\emptyset$ and (in the case $O\not\in V$) $V_c\cap W_c=\emptyset$. We need the following lemma.

\begin{lemma}[Cut by hyperplanes]\label{lem_hyperplane}
Let $V\subseteq\proj^m$ be a projective variety of dimension $r$ and degree $d$. Let $r'\geq r+1$. Choose $c_{i,j}\in S$ at random, for $i\in [r']$ and $0\leq j\leq m$.
Let $W\subseteq \proj^m$ be the projective subspace cut out by the equations $\sum_{j=0}^{m} c_{i,j} X_j=0$, $i=1,\dots,r'$, where $X_0,\dots,X_m$ are homogeneous coordinates of $\proj^m$. Then $V\cap W=\emptyset$ holds with probability at least $1-(r+1)d/|S|$.
\end{lemma}
\begin{proof}
For $i\in [r']$, let $H_i\subseteq \proj^m$ be the hyperplane defined by $\sum_{j=0}^{m} c_{i,j} X_j=0$. By ignoring $H_i$ for $i>r+1$, we may assume $r'=r+1$.  Let $V_0:= V$ and $V_i:= V_{i-1}\cap H_i$ for $i\in [r']$. It suffices to show that $\dim V_i=\dim V_{i-1}-1$ holds with probability at least $1-d/|S|$, for each $i\in [r']$ (the dimension of the empty set is $-1$ by convention). 

Fix $i\in [r']$ and $c_{i',j}$, for $i'\in [i-1]$ and $0\leq j\leq m$. So $V_{i-1}$ is also fixed.  Note that $V_{i-1}\neq \emptyset$ since  taking a hyperplane section reduces the dimension  by at most one. If $\dim V_i\neq \dim V_{i-1}-1$, then $\dim V_i=\dim V_{i-1}$, and $H_i$ contains some irreducible component of $V_{i-1}$~\cite[Exercise~11.6]{Har92}. Let $Y$ be an irreducible component of $V_{i-1}$, and fix a point $P\in Y$. Then $Y\subseteq H_i$ only if $P\in H_i$, which holds only if  $c_{i,0},\dots,c_{i,m}$ satisfy a nonzero linear equation determined by $P$. This occurs with probability at most $1/|S|$ (\eg, by fixing all but one $c_{i,j}$). 
We also have $\deg(V_{i-1})\leq \deg(V)\leq d$, and hence the number of irreducible components of $V_{i-1}$ is bounded by $d$.
By the union bound, $H_i$ contains an irreducible component of $V_{i-1}$ with probability at most $d/|S|$.
\end{proof}
%
%
\begin{proof}[Proof of \expref{Theorem}{thm-randReduct}]
As mentioned above, \expref{Theorem}{thm-randReduct} is equivalent to showing that, with probability at least $1-\delta$:  (1) $\dim V'=k$, and (2) $O'\in V'$ iff  $O\in V$.
Note that $W_c$ is cut out in $\proj^m$ by the linear equations $\sum_{j=1}^{m} c_{i,j} X_j=0$, $i=1,\dots,k+1$.
So $W_H$ is cut out in $H\cong \proj^{m-1}$ (corresponding to $X_0=0$) by the linear equations $\sum_{j=1}^m c_{i,j} X_j=0$, $i=1,\dots,k+1$. 
We also have $\deg(V_c\cap H)\leq \deg(V_c)\leq (\max_{i\in [m]}\deg(f_i))^k$ (see, \eg, \cite[Theorem~8.48]{burgisser2013algebraic}).

Assume $O\in V$. Then $O'\in V'$ since $\pi(O)=O'$. 
Applying \expref{Lemma}{lem_hyperplane} to  each of the irreducible components of $V_c\cap H$ and $W_H$, as subvarieties of $H\cong\proj^{m-1}$, we see $V_c\cap W_H=(V_c\cap H)\cap W_H=\emptyset$ holds with probability at least $1-k\deg(V_c \cap H)/|S|\geq 1-\delta$.
So by \expref{Lemma}{lem_projcriterion}, $\dim V'=k$ holds with probability at least $1-\delta$.

Now assume $O\not\in V$. 
Let $\pi_{O, H}: V_c\to H$ be the \emph{projection of $V_c$ from $O$ to $H$}, defined by $P\mapsto \overline{OP}\cap H$ for $P\in V_c$. It is well defined since $O\not\in V_c$. 
The image $\pi_{O,H}(V_c)$ is a projective subvariety of $H$~\cite[Theorem~3.5]{Har92}. 
If $V_c\cap W_c$ contains a point $P$, then $\pi_{O,H}(V_c)\cap W_H$ contains $\pi_{O,H}(P)$.
Conversely, if $\pi_{O,H}(V_c)\cap W_H$ contains a point $Q$, then there exists $P\in V_c$ such that $Q=\pi_{O,H}(P)$, and we have $P\in\overline{O Q}\subseteq W_c$.
We conclude that $\pi_{O,H}(V_c)\cap W_H=\emptyset$ iff  $V_c\cap W_c=\emptyset$, which implies  $V_c\cap W_H=\emptyset$. 

Note that $\dim \pi_{O,H}(V_c)=\dim V_c=k$, since
\[
  \pi_{O,H}(V_c)=J(\{O\}, V_c)\cap H\,.
\]
We also have $\deg(\pi_{O,H}(V_c))\leq \deg(V_c)$
\cite[Example~18.16]{Har92}.  Applying \expref{Lemma}{lem_hyperplane}
to $\pi_{O,H}(V_c)$ and $W_H$, as subvarieties of $H\cong\proj^{m-1}$,
we see $\pi_{O,H}(V_c)\cap W_H=\emptyset$ holds with probability at
least
\[
  1-\frac{(k+1)\deg(\pi_{O,H}(V_c))}{|S|} \geq 1-\delta\,.
\]

We have shown above that  $\pi_{O,H}(V_c)\cap W_H=\emptyset$ imply $V_c\cap W_H=\emptyset$ and $V_c\cap W_c=\emptyset$.
By  \expref{Lemma}{lem_projcriterion} and \expref{Lemma}{lem_projcriterion2},  $V_c\cap W_H=\emptyset$ and $V_c\cap W_c=\emptyset$ imply  $\dim V'=k$ and $O'\not\in V'$, respectively.
%
Therefore,  it holds with probability at least $1-\delta$ that $\dim V'=k$ and $O'\not\in V'$.
\end{proof}
%
%
\begin{proof}[Proof of \expref{Theorem}{thm-aps}]
$\AnnAtZero$ is known to be $\NP$-hard \cite{Kay09}. The $\NP$-hardness of $\APS$ follows from \expref{Lemma}{lem_geom} and \expref{Theorem}{thm_approx}.

Given an instance $\mathbf f$ of $\APS$, we can first find $k=\trdeg(\mathbf f)$. Fix a 
%
set %
$S\subseteq \aff$ to be larger than $2(k+1)(\max_{i\in [m]}\deg(f_i))^k$ (which can be scanned using only 
%
polynomial space). %
Consider the points $\left(\left(c_{i,j} \,\mid\, i\in [k+1],\, j\in [m] \right)\right)\in S^{(k+1)\times m}$; for each such point define $\mathbf g:= \big\{g_i:= \sum_{j=1}^m c_{i,j} f_j \,\mid$ $i\in[k+1] \big\}$.   
Compute the transcendence degree of $\mathbf g$, and if it is $k$ then solve $\AnnAtZero$ for the instance $\mathbf g$. 
Output NO iff some $\mathbf g$ failed the $\AnnAtZero$ test.

All these steps can be achieved in space polynomial in the input size, using the uniqueness of the annihilator for $\mathbf g$~\cite[Theorem~47]{Mat80}, Perron's degree bound~\cite{Plo05} and linear algebra~\cite{C76, BGH82, Mul87}.
\end{proof}
%
%
\section{Hitting set for \texorpdfstring{$\VPbar$}{VPbar}: Proof of \expref{Theorem}{thm-hsg}} \label{sec-hsg}

Suppose $p$ is a prime. Define $\aff:=\ol{\F}_p$. We want to find hitting sets for certain polynomials in $\aff[x_1,\dots,x_n]$. Fix a $p$-power $q\ge\Omega(sr^6)$, for the given parameters $s, r$. Assume that $p\nmid(r+1)$. Also, fix a model for the finite field $\F_q$~\cite{AL86}. We now define the notion of ``infinitesimally approximating'' a polynomial by a small circuit.

\smallskip\noindent
%
%
\textbf{Approximative closure of $\VP$.}~\cite{bringmann2017algebraic}
A family $(f_n|n)$ of polynomials from $\aff[\mathbf x]$ is in the \emph{class
%
%
$\VPbar_\aff$}       %
if there are polynomials $f_{n,i}$ and a function $t:\mathbb{N} \mapsto \mathbb{N}$ such that $g_n$  has a 
%
algebraic circuit of  size $\mathrm{poly}(n)$ and  degree $\mathrm{poly}(n)$, %
over the field $\aff(\epsilon)$, computing $g_n(\mathbf x)= f_n(\mathbf x)+ \epsilon f_{n,1}(\mathbf x)+ {\epsilon}^2 f_{n,2}(\mathbf x)+ \cdots+ {\epsilon}^{t(n)} f_{n,t(n)}(\mathbf x)$. That is, $g_n \equiv f_n \bmod{\epsilon\aff[\epsilon][\mathbf x]}$.
  
The smallest possible circuit size of $g_n$ is called the \emph{approximative complexity} of $f_n$, namely $\ol{\sz}(f_n)$. 

It may happen that $g_n$ is much easier than $f_n$ in terms of traditional circuit complexity. That possibility makes the definition interesting and opens up a long line of research. 

\smallskip\noindent
%
\textbf{Hitting set for $\VPbar_\aff$}       %
Given functions $s=s(n)$ and $r=r(n)$, a finite
%
set %
$\calH\subseteq \aff^n$ is called a \emph{hitting set} for degree-$r$ polynomials of approximative complexity $s$, if for every such nonzero polynomial $f$: $\exists \mathbf v\in\calH, \, f(\mathbf v)\ne0$.

\smallskip\noindent
\textbf{Explicitness.} Heintz and Schnorr~\cite{heintz1980testing} proved that 
%
$\mathrm{poly}(n,s)$-sized %
hitting sets exist aplenty (for degree-$r$ $\ol{\sz}$-$s$ %
polynomials) as follows. %

%
\begin{lemma}[{\cite[Theorem~4.4]{heintz1980testing}}]\label{lem-hs80}
 There exists a 
 %
 hitting set %
 $\mathcal{H}\subseteq \F_q^n $ of size $O(s^2n^2)$ (assuming $q\geq\Omega(sr^2)$) that hits all nonzero degree-$r$ $n$-variate polynomials in $\aff[\mathbf x]$ that can be infinitesimally approximated by size-$s$ algebraic circuits. 
\end{lemma}

The ultimate goal is computing such a hitting set in 
%
time $\mathrm{poly}(n,s,\log qr)$. %
Before our work, the best result known was
%
$\EXPSPACE$~\cite{mul12, mulmuley2017geometric}.

Note that for the hitting-set design problem it suffices to focus only on homogeneous polynomials. They are known to be computable by homogeneous circuits, where each gate computes a homogeneous polynomial~\cite{shpilka2010arithmetic}.

\smallskip\noindent
\textbf{Universal circuit.}
It can simulate any circuit of size $s$ computing a degree-$r$ homogeneous polynomial in $\aff(\epsilon)[x_1,\ldots,x_n]$. We define the \emph{universal circuit} $\Psi(\mathbf{y},\mathbf{x})$ as a circuit in $n$ essential variables $\mathbf{x}$ and $s':=O(sr^4)$ auxiliary variables $\mathbf y$. The variables $\mathbf y$ are the ones that one can specialize in $\aff(\epsilon)$, to compute a specific polynomial in $\aff(\epsilon)[x_1,\ldots,x_n]$. Every specialization gives a homogeneous degree-$r$ $\ol{\sz}$-$s'$ polynomial. Moreover, the set of these polynomials is closed under constant multiples~\cite[Theorem~2.2]{forbes2017pspace}.

Note that by~\cite{heintz1980testing} there is a hitting set, with $m:=O(s'^2n^2)$ points in $\F_q^n$ ($\because q\ge\Omega(s'r^2)$), 
%
for the set $\calP$ of the polynomials %
infinitesimally approximated by the specializations of $\Psi(\mathbf{y},\mathbf{x})$.
A universal circuit construction can be found in~\cite{raz2008elusive, shpilka2010arithmetic}. 
Using the above notation, we give a criterion to decide whether a candidate set is a hitting set. 

\begin{theorem}[hitting-set criterion]\label{thm_HITTING}
Set $\calH=:\{\mathbf v_1,\ldots, \mathbf v_m\}$ $\subseteq \F_q^n$ is \emph{not} a hitting set 
%
for the family $\calP$ of the polynomials %
infinitesimally approximated by the specializations of $\Psi(\mathbf{y},\mathbf{x})$ iff there is a satisfying assignment $(\alpha, \beta)\,\in \aff(\epsilon)^{s'}\times \aff(\epsilon)^{n}$ such that:

\smallskip\noindent
(1) $\forall i \in [n],\, {\beta_i}^{r+1}- 1 \,\in \epsilon\aff[\epsilon]$, where $r$ is the degree of the  specializations of $\Psi(\mathbf{y},\mathbf{x})$,  

\noindent 
(2) $\Psi(\alpha, \beta) - 1 \,\in \epsilon\aff[\epsilon]$, and

\noindent
(3) $\forall i \in [m],\, \Psi(\alpha,\mathbf{v}_i) \,\in \epsilon\aff[\epsilon]$.
\end{theorem}

%
\begin{remark}   %
The above criterion holds for algebraically closed fields $\aff$ of \emph{any} characteristic. Thus, it reduces 
%
the hitting set verification problems for the family $\calP$ over these fields %
to $\APS$ as well.
\end{remark}

\begin{proof}
%
First we show that %
$\exists x\in \aff(\epsilon),\, {x}^{r+1}- 1 \,\in \epsilon\aff[\epsilon]$ implies $x\in \aff[[\epsilon]]\cap \aff(\epsilon)$ (= rational functions defined at $\epsilon=0$).

\begin{claim}\label{defined at zero}
  $\exists x\in \aff(\epsilon),\, {x}^{r+1}- 1 \,\in \epsilon\aff[\epsilon]$ implies $x\in \, Z_{r+1} + \epsilon\aff[[\epsilon]]$, where 
  $Z_{r+1}$ is the set of $(r+1)$-th roots of unity in $\aff$. 
\end{claim}

 \begin{proof}

Recall the formal power series $\aff[[\epsilon]]$ and its group of units $\aff[[\epsilon]]^*$.
Note that for any polynomial
\[
  a=\left(\sum_{i_0\le i\le d} a_i\epsilon^i\right)
\]
with $a_{i_0}\ne0$, the inverse
\[
  a^{-1} = \epsilon^{-i_0}\cdot \left(\sum_{i_0\le i\le d}
  a_i\epsilon^{i-i_0}\right)^{-1}
\]
is in $\epsilon^{-i_0}\cdot \aff[[\epsilon]]^*$. This is just a
consequence of the identity
$(1-\epsilon)^{-1} = \sum_{i\ge0} \epsilon^i$. In other words, any
rational function $a\in \aff(\epsilon)$ can be written as an element
in $\epsilon^{-i}\aff[[\epsilon]]^*$, for some $i\ge0$. Thus, write
$x$ as $\epsilon^{-i}\cdot(b_0+b_1\epsilon+\cdots)$ for $i\ge0$ and
$b_0\in\aff^*$. This gives
$${x}^{r+1}- 1 \,=\, \epsilon^{-i(r+1)}(b_0+b_1\epsilon+ b_2\epsilon^{2} +\cdots)^{r+1} \;-\; 1 \,.$$
For this to be in $\epsilon\aff[\epsilon]$, clearly $i$ has to be $0$ (otherwise, $\epsilon^{-i(r+1)}$ remains 
%
uncancelled), %
implying that $x\in \aff[[\epsilon]]$. 

Moreover, we deduce that $b_0^{r+1} - 1 = 0$. Thus, 
%
Condition (1) %
implies that $b_0$ is one of the $(r+1)$-th roots of unity $Z_{r+1}\subseteq\aff$ (recall that, since $p\nmid(r+1)$, $|Z_{r+1}|=r+1$). Thus, $x\in \, Z_{r+1} + \epsilon\aff[[\epsilon]]$.
\end{proof}

\smallskip\emph{[$\Rightarrow$]:}
Suppose $\mathcal{H}$ is not a hitting set for $\calP$. Then, there is a specialization $\alpha \in \aff(\epsilon)^{s'}$ of the universal circuit such that $\Psi(\alpha, \mathbf x)$ computes a polynomial in $\aff[\epsilon][\mathbf x]\setminus \epsilon\aff[\epsilon][\mathbf x]$, but still ``fools'' $\calH$, \ie, $\forall i \in [m],\, \Psi(\alpha,\mathbf{v}_i) \,\in \epsilon\aff[\epsilon]$. 
What remains to be shown is that 
%
Conditions %
(1) and (2) 
can be satisfied too.

Consider the polynomial $g(\mathbf x):= \Psi(\alpha, \mathbf x)\big\vert_{\epsilon=0}$. It is a nonzero polynomial, in $\aff[\mathbf x]$ of 
%
degree $r$, %
that ``fools'' $\calH$. By 
%
\expref{Lemma}{Polynomial Identity}, %
there is a $\beta\in Z_{r+1}^n$ such that $a:= g(\beta)$ is in $\aff^*$. Clearly, $\beta_i^{r+1} - 1 = 0$, for all $i$. Consider $\psi':= a^{-1}\cdot\Psi(\alpha, \mathbf x)$. Note that $\psi'(\beta)-1 \in \epsilon\aff[\epsilon]$, and $\psi'(\mathbf{v}_i) \,\in \epsilon\aff[\epsilon]$ for all $i$. Moreover, the normalized polynomial $\psi'(\mathbf x)$ can easily be obtained from the universal circuit $\Psi$ by changing one of the coordinates of $\alpha$ (\eg, the incoming wires of the root of the circuit). This means that the three conditions 
%
(1) -- (3) %
can be simultaneously satisfied by (some) $(\alpha',\beta)\in\, \aff(\epsilon)^{s'}\times Z_{r+1}^n$.

\smallskip\emph{[$\Leftarrow$]:}
Suppose the satisfying assignment is $(\alpha, \beta')\,\in \aff(\epsilon)^{s'}\times \aff(\epsilon)^{n}$. As shown in \expref{Claim}{defined at zero}, 
%
Condition (1) implies %
$\beta'_i\in \, Z_{r+1} + \epsilon\aff[[\epsilon]]$ for all $i\in[n]$. Let us define $\beta_i:= \beta'_i\big\vert_{\epsilon=0}$, for all $i\in[n]$; they are in $Z_{r+1}\subseteq\aff$. 
%
By Condition (3), %
$\forall i \in [m],\, \Psi(\alpha,\mathbf{v}_i) \,\in \epsilon\aff[\epsilon]$.

Since any rational function $a\in \aff(\epsilon)$ can be written as an element in $\epsilon^{-i}\aff[[\epsilon]]^*$, for some $i\ge0$, we get that $\Psi(\alpha,\mathbf x)$ is in 
$\epsilon^{-j}\aff[[\epsilon]][\mathbf x]$, for some $j\ge0$.
Expand the polynomial $\Psi(\alpha,\mathbf x)$, w.r.t.\ $\epsilon$, 
%
as %
$$g_{-j}(\mathbf x)\epsilon^{-j} + \dots + \epsilon^{-2} g_{-2}(\mathbf x) + g_{-1}(\mathbf x)\epsilon^{-1} + g_0(\mathbf x) + \epsilon g_1(\mathbf x) + \epsilon^2 g_2(\mathbf x) + \dots \,.$$
%
Let us study Condition (2). If for each $0\le\ell\le j$, polynomial $g_{-\ell}(\mathbf x)$ is zero, then $\Psi(\alpha,\beta')\big\vert_{\epsilon=0} =0$ contradicting the condition. Thus, we can pick the largest $0\le\ell\le j$ such that the polynomial $g_{-\ell}(\mathbf x)\ne0$. %
%

Note that the normalized circuit $\epsilon^\ell\cdot \Psi(\alpha,\mathbf x)$ equals $g_{-\ell}$ at $\epsilon=0$. This means that $g_{-\ell}\in\calP$, and it is a nonzero polynomial fooling $\calH$. Thus, $\calH$ cannot be a hitting set for $\calP$ and we are done.
\end{proof}

\begin{proof}[Proof of \expref{Theorem}{thm-hsg}]
Given a prime $p$ and parameters $n, r, s$ in unary (w.l.o.g.\ $p\nmid(r+1)$), fix a field $\F_q$ with $q\ge\Omega(sr^6)$. Fix the universal circuit $\Psi(\mathbf{y},\mathbf{x})$ with $n$ essential variables $\mathbf{x}$ and $s':=\Omega(sr^4)$ auxiliary variables $\mathbf y$. Fix $m:=\Omega(s'^2n^2)$.

For every subset $\calH=:\{\mathbf v_1,\ldots, \mathbf v_m\}$ $\subseteq \F_q^n$ solve the $\APS$ instance described by 
Conditions 
%
(1) -- (3) %
in \expref{Theorem}{thm_HITTING}. These are $(n+m+1)$ algebraic circuits of degree 
%
$\mathrm{poly}(srn, \log p)$ and a similar bit size. %
Using the algorithm from \expref{Theorem}{thm-aps} it can be solved in 
%
space $\mathrm{poly}(srn, \log p)$. %

The number of subsets $\calH$ is 
at most %
$q^{nm}$. 
%
So we can go over all of them in space $\mathrm{poly}(nm\log q)$. %
If $\APS$ fails on one of them (say $\calH$) then we know that $\calH$ is a hitting set for $\calP$. Since $\Psi$ is universal, for homogeneous degree-$r$ $\ol{\sz}$-$s$ polynomials in $\aff[\mathbf x]$, we output $\calH$ as the desired hitting set.
\end{proof}

%
\begin{remark}   %
One advantage of our method compared to the one in~\cite{forbes2017pspace} is that
%
in our construction, the 
%
bit complexity of the coordinates of the points in the hitting set %
is $O(\log rs)$, whereas the
%
bit complexity of those in~\cite{forbes2017pspace} is  $\mathrm{poly}(n,s,r)$. %
\end{remark}

%

\section{Conclusion}\label{sec-conclusion}
%

Our result
%
that  %
algebraic dependence testing
is  %
in
%
$\AMcapcoAM$   %
gives
%
%
some hope   %
that a randomized
%
polynomial-time   %
algorithm for the problem
%
might exist.  %
Studying the following special case might be helpful to get an idea for designing better algorithms.

  Given quadratic polynomials $f_1,\ldots,f_n \in \F_2[x_1,\ldots,x_n]$, test if they are algebraically dependent in randomized polynomial time~\cite{pandey2016algebraic}.
  
 
In addition, Ilya Volkovich~\cite{Vol18} asked
%
%
%
%
whether $\AD(\F) \in \SBPcapcoSBP$ is true, and whether our method can be used to prove this stronger result. %
%
%
%
%
Here
%
%
%
%
$\SBP$, introduced in~\cite{BCM06},  %
stands for ``small bounded-error probability'' and is a subclass of $\AM$. %


\smallskip
 As indicated in this paper, approximate polynomials satisfiability, or equivalently testing zero-membership in the Zariski closure of the image, may have further applications to problems in computational algebraic geometry and algebraic complexity. 

 We know that
 %
 $\HN$ %
 is in $\AM$ over fields of characteristic zero, assuming GRH~\cite{koiran1996hilbert}. Can we 
%
prove that $\AnnAtZero$ (or $\APS$) is also in $\AM$ over %
 fields of characteristic zero assuming GRH~\cite{Kay09}? This would also imply a better hitting set construction for %
 $\VPbar$.


\bigskip\noindent
\textbf{Acknowledgements.}
We thank Anurag Pandey and Sumanta Ghosh for insightful discussions on the approximate polynomials satisfiability and explicit construction of hitting set for the closure of 
%
$\VP$. %
We are also grateful to the anonymous reviewers for their careful reading of the manuscript and for their many helpful comments. We thank the editor-in-chief Laci Babai for his meticulous suggestions regarding grammar and style.
Finally, N.S. thanks the funding support from DST (DST/SJF/MSA-01/2013-14), and Z.G. thanks the funding support from DST and Research I Foundation of CSE, IITK.

\nocite{lidl}

%% Theory of Computing
%% File: v015a016.bbl
%% Publication date: December 14, 2019
%% Authors: Zeyu Guo, Nitin Saxena, and Amit Sinhababu
% Define empty bibhead if not already defined
\providecommand{\bibhead}[1]{}
% Define tocrefpdfbookmark if not already defined
\expandafter\ifx\csname pdfbookmark\endcsname\relax%
  \providecommand{\tocrefpdfbookmark}{}
\else\providecommand{\tocrefpdfbookmark}{%
   \hypertarget{tocreferences}{}%
   \pdfbookmark[1]{References}{tocreferences}}%
\fi

\tocrefpdfbookmark
\begin{thebibliography}{10}

\bibitem{AL86}\bibhead{AL86}
{\sc Leonard~M. Adleman and Hendrik~W. Lenstra}: Finding irreducible
  polynomials over finite fields.
\newblock In {\em Proc. 18th STOC}, pp. 350--355. ACM Press, 1986.
\newblock [\epfmtdoi{10.1145/12130.12166}]

\bibitem{AGS17}\bibhead{AGS17}
{\sc Manindra Agrawal, Sumanta Ghosh, and Nitin Saxena}: Bootstrapping
  variables in algebraic circuits.
\newblock {\em Proc. Nat. Acad. of Sciences (USA)}, 116(17):8107--8118, 2019.
\newblock Preliminary version in
  \href{https://dl.acm.org/citation.cfm?id=3188762}{STOC'18}.
\newblock [\epfmtdoi{10.1073/pnas.1901272116}]

\bibitem{ASSS11}\bibhead{ASSS11}
{\sc Manindra Agrawal, Chandan Saha, Ramprasad Saptharishi, and Nitin Saxena}:
  Jacobian hits circuits: {H}itting sets, lower bounds for depth-{$D$}
  occur-$k$ formulas and depth-3 transcendence degree-$k$ circuits.
\newblock {\em SIAM J. Comput.}, 45(4):1533--1562, 2016.
\newblock Preliminary version in
  \href{https://dl.acm.org/citation.cfm?id=2214033&dl=ACM&coll=DL}{STOC'12}.
\newblock [\epfmtdoi{10.1137/130910725}, \epfmt{arxiv}{1111.0582}]

\bibitem{AB09}\bibhead{AB09}
{\sc Sanjeev Arora and Boaz Barak}: {\em Computational Complexity: A Modern
  Approach}.
\newblock Cambridge Univ. Press, 2009.
\newblock [\epfmtdoi{10.1017/CBO9780511804090}]

\bibitem{arvind}\bibhead{arvind}
{\sc Vikraman Arvind, Pushkar~S. Joglekar, Partha Mukhopadhyay, and S.~Raja}:
  Randomized polynomial-time identity testing for noncommutative circuits.
\newblock {\em Theory of Computing}, 15(7):1--36, 2019.
\newblock Preliminary version in
  \href{https://doi.org/10.1145/3055399.3055442}{STOC'17}.
\newblock [\epfmtdoi{10.4086/toc.2019.v015a007}]

\bibitem{babai1985trading}\bibhead{babai1985trading}
{\sc L{\'a}szl{\'o} Babai}: Trading group theory for randomness.
\newblock In {\em Proc. 17th STOC}, pp. 421--429. ACM Press, 1985.
\newblock [\epfmtdoi{10.1145/22145.22192}]

\bibitem{BMS13}\bibhead{BMS13}
{\sc Malte Beecken, Johannes Mittmann, and Nitin Saxena}: Algebraic
  independence and blackbox identity testing.
\newblock {\em Inform. and Comput.}, 222:2--19, 2013.
\newblock Preliminary version in
  \href{https://link.springer.com/chapter/10.1007/978-3-642-22012-8_10}{ICALP'11}.
\newblock [\epfmtdoi{10.1016/j.ic.2012.10.004}, \epfmt{arxiv}{1102.2789}]

\bibitem{BCM06}\bibhead{BCM06}
{\sc Elmar B{\"o}hler, Christian Gla{\ss}er, and Daniel Meister}: Error-bounded
  probabilistic computations between {MA} and {AM}.
\newblock {\em J. Comput. System Sci.}, 72(6):1043--1076, 2006.
\newblock Preliminary version in
  \href{https://link.springer.com/chapter/10.1007/978-3-540-45138-9_19}{MFCS'03}.
\newblock [\epfmtdoi{10.1016/j.jcss.2006.05.001}]

\bibitem{BGH82}\bibhead{BGH82}
{\sc Allan Borodin, Joachim von~zur Gathem, and John Hopcroft}: Fast parallel
  matrix and {GCD} computations.
\newblock {\em Inf. Control}, 52(3):241--256, 1982.
\newblock Preliminary version in
  \href{https://ieeexplore.ieee.org/document/4568376}{FOCS'82}.
\newblock [\epfmtdoi{10.1016/S0019-9958(82)90766-5}]

\bibitem{bringmann2017algebraic}\bibhead{bringmann2017algebraic}
{\sc Karl Bringmann, Christian Ikenmeyer, and Jeroen Zuiddam}: On algebraic
  branching programs of small width.
\newblock {\em J. ACM}, 65(5):32:1--32:29, 2018.
\newblock Preliminary version in
  \href{http://drops.dagstuhl.de/opus/volltexte/2017/7521/}{CCC'17}.
\newblock [\epfmtdoi{10.1145/3209663}, \epfmt{arxiv}{1702.05328}]

\bibitem{burgisser2004complexity}\bibhead{burgisser2004complexity}
{\sc Peter B{\"u}rgisser}: The complexity of factors of multivariate
  polynomials.
\newblock {\em Foundations of Computational Mathematics}, 4(4):369--396, 2004.
\newblock Preliminary version in
  \href{https://doi.org/10.1109/SFCS.2001.959912}{FOCS'01}.
\newblock [\epfmtdoi{10.1007/s10208-002-0059-5}]

\bibitem{burgisser2013algebraic}\bibhead{burgisser2013algebraic}
{\sc Peter B{\"u}rgisser, Michael Clausen, and Amin Shokrollahi}: {\em
  Algebraic Complexity Theory}.
\newblock Springer, 1997.
\newblock [\epfmtdoi{10.1007/978-3-662-03338-8}]

\bibitem{brgisser_et_al:LIPIcs:2018:8351}\bibhead{brgisser_et_al:LIPIcs:2018:8351}
{\sc Peter B{\"u}rgisser, Ankit Garg, Rafael Oliveira, Michael Walter, and Avi
  Wigderson}: Alternating minimization, scaling algorithms, and the null-cone
  problem from invariant theory.
\newblock In {\em Proc. 9th Innovations in Theoretical Computer Science Conf.
  (ITCS'18)}, pp. 24:1--24:20. Schloss Dagstuhl--Leibniz-Zentrum fuer
  Informatik, 2018.
\newblock [\epfmtdoi{10.4230/LIPIcs.ITCS.2018.24}, \epfmt{arxiv}{1711.08039}]

\bibitem{CKS}\bibhead{CKS}
{\sc Chi-Ning Chou, Mrinal Kumar, and Noam Solomon}: Closure results for
  polynomial factorization.
\newblock {\em Theory of Computing}, 15(13):1--34, 2019.
\newblock Preliminary version in
  \href{https://doi.org/10.4230/LIPIcs.CCC.2018.13}{CCC'18}.
\newblock [\epfmtdoi{10.4086/toc.2019.v015a013}]

\bibitem{C76}\bibhead{C76}
{\sc Laszlo Csanky}: Fast parallel matrix inversion algorithms.
\newblock {\em SIAM J. Comput.}, 5(4):618--623, 1976.
\newblock [\epfmtdoi{10.1137/0205040}]

\bibitem{demillo1978probabilistic}\bibhead{demillo1978probabilistic}
{\sc Richard~A. Demillo and Richard~J. Lipton}: A probabilistic remark on
  algebraic program testing.
\newblock {\em Inform. Process. Lett.}, 7(4):193--195, 1978.
\newblock [\epfmtdoi{10.1016/0020-0190(78)90067-4}]

\bibitem{derksen2015computational}\bibhead{derksen2015computational}
{\sc Harm Derksen and Gregor Kemper}: {\em Computational Invariant Theory}.
\newblock Springer, 2015.
\newblock [\epfmtdoi{10.1007/978-3-662-48422-7}]

\bibitem{bib:Dvi09}\bibhead{bib:Dvi09}
{\sc Zeev Dvir}: Extractors for varieties.
\newblock {\em Comput. Complexity}, 21(4):515--572, 2012.
\newblock Preliminary version in
  \href{https://ieeexplore.ieee.org/document/5231245}{CCC'09}.
\newblock [\epfmtdoi{10.1007/s00037-011-0023-3}]

\bibitem{DGW09}\bibhead{DGW09}
{\sc Zeev Dvir, Ariel Gabizon, and Avi Wigderson}: Extractors and rank
  extractors for polynomial sources.
\newblock {\em Comput. Complexity}, 18(1):1--58, 2009.
\newblock Preliminary version in
  \href{https://ieeexplore.ieee.org/document/4389479}{FOCS'07}.
\newblock [\epfmtdoi{10.1007/s00037-009-0258-4}]

\bibitem{ehrenborg1993apolarity}\bibhead{ehrenborg1993apolarity}
{\sc Richard Ehrenborg and Gian-Carlo Rota}: Apolarity and canonical forms for
  homogeneous polynomials.
\newblock {\em Eur. J. Combinatorics}, 14(3):157--181, 1993.
\newblock [\epfmtdoi{10.1006/eujc.1993.1022}]

\bibitem{forbes2017pspace}\bibhead{forbes2017pspace}
{\sc Michael~A. Forbes and Amir Shpilka}: A {PSPACE} construction of a hitting
  set for the closure of small algebraic circuits.
\newblock In {\em Proc. 50th STOC}, pp. 1180--1192. ACM Press, 2018.
\newblock [\epfmtdoi{10.1145/3188745.3188792}]

\bibitem{goldwasser1986private}\bibhead{goldwasser1986private}
{\sc Shafi Goldwasser and Michael Sipser}: Private coins versus public coins in
  interactive proof systems.
\newblock In {\sc Silvio Micali}, editor, {\em Randomness in Computation},
  volume~5 of {\em Advances in Computing Research}, pp. 73--90. JAI Press,
  1989.
\newblock Preliminary version in
  \href{https://doi.org/10.1145/12130.12137}{STOC'86}.

\bibitem{grochow2018circuit}\bibhead{grochow2018circuit}
{\sc Joshua~A. Grochow and Toniann Pitassi}: Circuit complexity, proof
  complexity, and polynomial identity testing: {T}he ideal proof system.
\newblock {\em J. ACM}, 65(6):37:1--37:59, 2018.
\newblock Preliminary version in
  \href{https://ieeexplore.ieee.org/document/6978995}{FOCS'14}.
\newblock [\epfmtdoi{10.1145/3230742}]

\bibitem{guo2018algebraic}\bibhead{guo2018algebraic}
{\sc Zeyu Guo, Nitin Saxena, and Amit Sinhababu}: Algebraic dependencies and
  {PSPACE} algorithms in approximative complexity.
\newblock In {\em Proc. 33rd Computational Complexity Conf. (CCC'18)}, pp.
  10:1--10:21. Schloss Dagstuhl--Leibniz-Zentrum fuer Informatik, 2018.
\newblock [\epfmtdoi{10.4230/LIPIcs.CCC.2018.10}, \epfmt{arxiv}{1801.09275}]

\bibitem{Har92}\bibhead{Har92}
{\sc Joe Harris}: {\em Algebraic Geometry: A First Course}.
\newblock Springer, 1992.
\newblock [\epfmtdoi{10.1007/978-1-4757-2189-8}]

\bibitem{hartshorne2013algebraic}\bibhead{hartshorne2013algebraic}
{\sc Robin Hartshorne}: {\em Algebraic Geometry}.
\newblock Springer, 1977.
\newblock [\epfmtdoi{10.1007/978-1-4757-3849-0}]

\bibitem{heintz1980testing}\bibhead{heintz1980testing}
{\sc Joos Heintz and Claus-Peter Schnorr}: Testing polynomials which are easy
  to compute.
\newblock In {\em Proc. 12th STOC}, pp. 262--272. ACM Press, 1980.
\newblock [\epfmtdoi{10.1145/800141.804674}]

\bibitem{ingleton1971representation}\bibhead{ingleton1971representation}
{\sc Aubrey~W. Ingleton}: Representation of matroids.
\newblock {\em Combinatorial Mathematics and its applications}, pp. 149--167,
  1971.

\bibitem{Jac41}\bibhead{Jac41}
{\sc Carl Gustav~Jacob Jacobi}: De {D}eterminantibus functionalibus.
\newblock {\em Journal f{\"u}r die reine und angewandte Mathematik (Crelles
  J.)}, 1841(22):319--359, 1841.
\newblock [\epfmtdoi{10.1515/crll.1841.22.319}]

\bibitem{Kay09}\bibhead{Kay09}
{\sc Neeraj Kayal}: The complexity of the annihilating polynomial.
\newblock In {\em Proc. 24th IEEE Conf. on Computational Complexity (CCC'09)},
  pp. 184--193. IEEE Comp. Soc. Press, 2009.
\newblock [\epfmtdoi{10.1109/CCC.2009.37}]

\bibitem{kayal2006complexity}\bibhead{kayal2006complexity}
{\sc Neeraj Kayal and Nitin Saxena}: Complexity of ring morphism problems.
\newblock {\em Comput. Complexity}, 15(4):342--390, 2006.
\newblock [\epfmtdoi{10.1007/s00037-007-0219-8}]

\bibitem{koiran1996hilbert}\bibhead{koiran1996hilbert}
{\sc Pascal Koiran}: Hilbert's {N}ullstellensatz is in the polynomial
  hierarchy.
\newblock {\em J. Complexity}, 12(4):273--286, 1996.
\newblock [\epfmtdoi{10.1006/jcom.1996.0019}]

\bibitem{kollar1988sharp}\bibhead{kollar1988sharp}
{\sc J{\'a}nos Koll{\'a}r}: Sharp effective {N}ullstellensatz.
\newblock {\em J. Amer. Math. Soc.}, 1(4):963--975, 1988.
\newblock [\epfmtdoi{10.2307/1990996}]

\bibitem{KS15}\bibhead{KS15}
{\sc Mrinal Kumar and Shubhangi Saraf}: Arithmetic circuits with locally low
  algebraic rank.
\newblock {\em Theory of Computing}, 13(6):1--33, 2017.
\newblock Preliminary version in
  \href{http://drops.dagstuhl.de/opus/volltexte/2016/5828/}{CCC'16}.
\newblock [\epfmtdoi{10.4086/toc.2017.v013a006}, \epfmt{arxiv}{1806.06097}]

\bibitem{landsberg2012tensors}\bibhead{landsberg2012tensors}
{\sc Joseph~M. Landsberg}: {\em Tensors: Geometry and Applications}.
\newblock Amer. Math. Soc., 2012.
\newblock [\epfmtdoi{10.1090/gsm/128}]

\bibitem{le2014powers}\bibhead{le2014powers}
{\sc Fran{\c{c}}ois Le~Gall}: Powers of tensors and fast matrix multiplication.
\newblock In {\em Proc. 39th Internat. Symp. Symbolic and Algebraic Computation
  (ISSAC'14)}, pp. 296--303. ACM Press, 2014.
\newblock [\epfmtdoi{10.1145/2608628.2608664}, \epfmt{arxiv}{1401.7714}]

\bibitem{LL89}\bibhead{LL89}
{\sc Thomas Lehmkuhl and Thomas Lickteig}: On the order of approximation in
  approximative triadic decompositions of tensors.
\newblock {\em Theoret. Comput. Sci.}, 66(1):1--14, 1989.
\newblock [\epfmtdoi{10.1016/0304-3975(89)90141-2}]

\bibitem{lidl}\bibhead{lidl}
{\sc Rudolf Lidl and Harald Niederreiter}: {\em Finite Fields}.
\newblock Encyclopedia of Mathematics and its Applications. Cambridge Univ.
  Press, 2nd edition, 1996.
\newblock [\epfmtdoi{10.1017/CBO9780511525926}]

\bibitem{Mat80}\bibhead{Mat80}
{\sc Hideyuki Matsumura}: {\em Commutative Algebra}.
\newblock Benjamin-Cummings Pub Co, 1980.

\bibitem{mayr1982complexity}\bibhead{mayr1982complexity}
{\sc Ernst~W. Mayr and Albert~R. Meyer}: The complexity of the word problems
  for commutative semigroups and polynomial ideals.
\newblock {\em Adv. Math.}, 46(3):305--329, 1982.
\newblock [\epfmtdoi{10.1016/0001-8708(82)90048-2}]

\bibitem{mittmann2014algebraic}\bibhead{mittmann2014algebraic}
{\sc Johannes Mittmann, Nitin Saxena, and Peter Scheiblechner}: Algebraic
  independence in positive characteristic: {A} $p$-adic calculus.
\newblock {\em Trans. Amer. Math. Soc.}, 366(7):3425--3450, 2014.
\newblock [\epfmtdoi{10.1090/S0002-9947-2014-06268-5},
  \epfmt{arxiv}{1202.4301}]

\bibitem{muller}\bibhead{muller}
{\sc David~E. Muller}: Application of boolean algebra to switching circuit
  design and to error detection.
\newblock {\em Trans. Inst. Radio Engineers Professional Group on Electronic
  Computers}, EC-3(3):6--12, 1954.
\newblock [\epfmtdoi{10.1109/IREPGELC.1954.6499441}]

\bibitem{Mul87}\bibhead{Mul87}
{\sc Ketan~D. Mulmuley}: A fast parallel algorithm to compute the rank of a
  matrix over an arbitrary field.
\newblock {\em Combinatorica}, 7(1):101--104, 1987.
\newblock Preliminary version in
  \href{https://dl.acm.org/citation.cfm?id=12164}{STOC'86}.
\newblock [\epfmtdoi{10.1007/BF02579205}]

\bibitem{mul12}\bibhead{mul12}
{\sc Ketan~D. Mulmuley}: Geometric complexity theory {V}: {E}quivalence between
  blackbox derandomization of polynomial identity testing and derandomization
  of {N}oether's normalization lemma.
\newblock In {\em Proc. 53rd FOCS}, pp. 629--638. IEEE Comp. Soc. Press, 2012.
\newblock [\epfmtdoi{10.1109/FOCS.2012.15}]

\bibitem{mulmuley2017geometric}\bibhead{mulmuley2017geometric}
{\sc Ketan~D. Mulmuley}: Geometric complexity theory {V}: {E}fficient
  algorithms for {Noether} normalization.
\newblock {\em J. Amer. Math. Soc.}, 30(1):225--309, 2017.
\newblock Preliminary version in
  \href{https://doi.org/10.1109/FOCS.2012.15}{FOCS'12}.
\newblock [\epfmtdoi{10.1090/jams/864}, \epfmt{arxiv}{1209.5993}]

\bibitem{ore}\bibhead{ore}
{\sc {\O}ystein Ore}: {\"U}ber h{\"o}here {K}ongruenzen.
\newblock {\em Norsk Mat. Forenings Skrifter Ser. I}, 7(15):27, 1922.
\newblock Polynomial Identity Lemma cited with full proof in \cite[Theorem
  6.13]{lidl}.

\bibitem{pandey2016algebraic}\bibhead{pandey2016algebraic}
{\sc Anurag Pandey, Nitin Saxena, and Amit Sinhababu}: Algebraic independence
  over positive characteristic: {N}ew criterion and applications to locally
  low-algebraic-rank circuits.
\newblock {\em Comput. Complexity}, 27(4):617--670, 2018.
\newblock Preliminary version in
  \href{https://doi.org/10.4230/LIPIcs.MFCS.2016.74}{MFCS'16}.
\newblock [\epfmtdoi{10.1007/s00037-018-0167-5}]

\bibitem{Per27}\bibhead{Per27}
{\sc Oskar Perron}: {\em Algebra. I: Die Grundlagen}.
\newblock Walter de Gruyter, 1927.

\bibitem{Plo05}\bibhead{Plo05}
{\sc Arkadiusz P{\l}oski}: Algebraic dependence of polynomials after {O.
  Perron} and some applications.
\newblock In {\em Computational Commutative and Non-Commutative Algebraic
  Geometry}, pp. 167--173. IOS Press, 2005.

\bibitem{Rai12}\bibhead{Rai12}
{\sc Claudiu Raicu}: Secant varieties of {Segre--Veronese} varieties.
\newblock {\em Algebra \& Number Theory}, 6(8):1817--1868, 2012.
\newblock [\epfmtdoi{10.2140/ant.2012.6.1817}]

\bibitem{raz2008elusive}\bibhead{raz2008elusive}
{\sc Ran Raz}: Elusive functions and lower bounds for arithmetic circuits.
\newblock {\em Theory of Computing}, 6(7):135--177, 2010.
\newblock Preliminary version in
  \href{https://dl.acm.org/citation.cfm?doid=1374376.1374479}{STOC'08}.
\newblock [\epfmtdoi{10.4086/toc.2010.v006a007}]

\bibitem{Saxena09}\bibhead{Saxena09}
{\sc Nitin Saxena}: Progress on polynomial identity testing.
\newblock {\em Bulletin of the {EATCS}}, 99:49--79, 2009.
\newblock Online version
  \href{https://eccc.weizmann.ac.il/report/2009/101/}{ECCC TR09-101}.

\bibitem{Saxena14}\bibhead{Saxena14}
{\sc Nitin Saxena}: Progress on polynomial identity testing - {II}.
\newblock In {\em Perspectives in Computational Complexity}, pp. 131--146.
  Springer, 2014.
\newblock [\epfmtdoi{10.1007/978-3-319-05446-9\_7}, \epfmt{arxiv}{1401.0976}]

\bibitem{schaefer2016complexity}\bibhead{schaefer2016complexity}
{\sc Marcus Schaefer and Daniel {\v{S}}tefankovi{\v{c}}}: The complexity of
  tensor rank.
\newblock {\em Theory of Computing Systems}, 62(5):1161--1174, 2018.
\newblock [\epfmtdoi{10.1007/s00224-017-9800-y}, \epfmt{arxiv}{1612.04338}]

\bibitem{Sch95}\bibhead{Sch95}
{\sc Joachim Schmid}: On the affine {B}ezout inequality.
\newblock {\em manuscripta mathematica}, 88(1):225--232, 1995.
\newblock [\epfmtdoi{10.1007/BF02567819}]

\bibitem{schmidt}\bibhead{schmidt}
{\sc Wolfgang~M. Schmidt}: {\em Equations over Finite Fields: An Elementary
  Approach}.
\newblock Volume 536 of {\em Lecture Notes in Math.}
\newblock Springer, 1st edition, 1976.
\newblock [\epfmtdoi{10.1007/BFb0080437}]

\bibitem{Sch80}\bibhead{Sch80}
{\sc Jacob~T. Schwartz}: Fast probabilistic algorithms for verification of
  polynomial identities.
\newblock {\em J. ACM}, 27(4):701--717, 1980.
\newblock Preliminary version in
  \href{https://link.springer.com/chapter/10.1007/3-540-09519-5_72}{EUROSAM'79}.
\newblock [\epfmtdoi{10.1145/322217.322225}]

\bibitem{shpilka2010arithmetic}\bibhead{shpilka2010arithmetic}
{\sc Amir Shpilka and Amir Yehudayoff}: Arithmetic circuits: {A} survey of
  recent results and open questions.
\newblock {\em Foundations and Trends in Theoretical Computer Science},
  5(3--4):207--388, Now Publ., 2009.
\newblock [\epfmtdoi{10.1561/0400000039}]

\bibitem{Vol18}\bibhead{Vol18}
{\sc Ilya Volkovich}: Private Communication, 2018.

\bibitem{zippel1979probabilistic}\bibhead{zippel1979probabilistic}
{\sc Richard~E. Zippel}: Probabilistic algorithms for sparse polynomials.
\newblock In {\em Proc. Internat. Symp. Symbolic and Algebraic Manipulation
  (EUROSAM'79)}, pp. 216--226. Springer, 1979.
\newblock [\epfmtdoi{10.1007/3-540-09519-5\_73}]

\end{thebibliography}

\end{document}
